\chapter{Determinantal image and composition toward Birch--Swinnerton--Dyer}

\section{Genealogical unit and fiber product}

Let $E/\mathbb Q$ be an elliptic curve and let
$\mathfrak G_{\rm det}(E)$ be the complete autonomous determinantal image. Its
category of surviving histories is denoted by
$\mathcal T_m^{\rm surv}$. The preceding determinantal owner constructs from a
common unit the chain of lines
\eqref{eq:amp-bsd-linea-determinante}--\eqref{eq:amp-bsd-cadena-lineas}, the
Kato--FERS comparator
\eqref{eq:amp-bsd-phi-kato-fers}--\eqref{eq:amp-bsd-inversa-kato-fers}, the
Poitou--Tate channel
\eqref{eq:amp-bsd-comparador-pt-q}--\eqref{eq:amp-bsd-comparador-pt} and the
local--global gluing
\eqref{eq:amp-bsd-triangulo-local}--\eqref{eq:amp-bsd-mapa-localizacion}.
The tuple of \eqref{eq:pdf2-g-det-explicito} and the identities of
\eqref{eq:pdf2-g-det-identidades-dato} are the joint publication of that
upstream construction: they are not accepted as terminal axioms or chosen
from the rank or leading term. This chapter preserves and composes
those arrows through to the Birch--Swinnerton--Dyer conclusion.

The coefficient object is not merely the $3$-primary fibre. One first fixes
the based integral free complex
\begin{equation}\label{eq:pdf2-bsd-gmz}
 G_{m,\mathbb Z}:=N_*\bigl(N\mathcal T_m^{\rm surv};\mathbb Z\bigr),
\end{equation}
and, for each prime $\ell$ and each $n\geq1$, one sets
\begin{equation}\label{eq:pdf2-bsd-coeficientes-todos-ell}
 A_{\ell,n}:=E[\ell^n](\overline{\mathbb Q}),\qquad
 T_\ell E:=\varprojlim_nA_{\ell,n},\qquad
 V_\ell E:=T_\ell E\otimes_{\mathbb Z_\ell}\mathbb Q_\ell,
 \qquad
 G_{m,\ell,n}:=G_{m,\mathbb Z}\otimes_{\mathbb Z}^{\mathbf L}A_{\ell,n}.
\end{equation}
The transitions are not implicit: they are
\begin{equation}\label{eq:pdf2-bsd-transiciones-ell-n}
 q_{\ell,n}:A_{\ell,n+1}\twoheadrightarrow A_{\ell,n},
 \quad x\longmapsto\ell x,
 \qquad
 Q_{\ell,n}:=1\otimes q_{\ell,n}:G_{m,\ell,n+1}\to G_{m,\ell,n},
\end{equation}
The transition of coefficients in the Manin complex, together with
$Q_{\ell,n}$, induces
$\widetilde Q_{\ell,n}:P_{E,\ell,n+1}^{\rm gen}\to
P_{E,\ell,n}^{\rm gen}$. These maps satisfy
$Q_{\ell,n}d=dQ_{\ell,n}$,
$(Q_{\ell,n}\otimes Q_{\ell,n})\Delta_{\rm AW}
=\Delta_{\rm AW}Q_{\ell,n}$, and
$\widetilde Q_{\ell,n}\widehat\Pi_{0,\ell,n+1}
=\widehat\Pi_{0,\ell,n}Q_{\ell,n}$. Thus, the limit defining $T_\ell E$
is taken within the same family of complexes, not only in the coefficient
groups.
The former notation $A_n$ denotes only
$A_{3,n}$; it does not govern the arithmetic publication. The boundary matrices,
Alexander--Whitney, orientation, and carry are written over $\mathbb Z$ in
the basis of histories. Therefore, for
$R\in\{\mathbb Z_\ell,\mathbb Z/\ell^n,\mathbb Q_\ell\}$,
\begin{equation}\label{eq:pdf2-bsd-base-change-integral}
 G_{m,R}=G_{m,\mathbb Z}\otimes_{\mathbb Z}^{\mathbf L}R,
 \qquad
 d_{m,R}=d_{m,\mathbb Z}\otimes1,
 \qquad
 \Delta_{{\rm AW},R}=\Delta_{{\rm AW},\mathbb Z}\otimes1.
\end{equation}
Thus, passing from $3$ to another prime is not a nonexistent extension
$\mathbb Z_3\to\mathbb Z_\ell$: all fibers are obtained by base
change from the same integral owner.

For brevity, the following formulas are written over a fixed coefficient
$A_{\ell,n}$ and are compatible with $n$, with $\ell$, and with the
integral basis. The normalized complex
\begin{equation}\label{eq:pdf2-bsd-gmn}
 G_{m,\ell,n}:=N_*\bigl(N\mathcal T_m^{\rm surv};A_{\ell,n}\bigr)
\end{equation}
preserves the Alexander--Whitney diagonal
\begin{equation}\label{eq:pdf2-bsd-aw}
 \Delta_{\rm AW}[x_0\to\cdots\to x_q]
 =\sum_{j=0}^q[x_0\to\cdots\to x_j]\otimes
                [x_j\to\cdots\to x_q]
\end{equation}
and the counit
$\varepsilon_{\rm AW}[x]=1$ at vertices, zero in positive degrees. The history
couples to the Manin complex through
\begin{equation}\label{eq:pdf2-bsd-pullback}
 P_{E,\ell,n}^{\rm gen}:=
 P_{E,\ell,n}^{\rm Man}\times_{A_{\ell,n}[0]}G_{m,\ell,n}.
\end{equation}
If $s_{\ell,n}:A_{\ell,n}[0]\to P_{E,\ell,n}^{\rm Man}$ is the based section, then
\[
 \widehat\Pi_{0,\ell,n}(c)=
 (s_{\ell,n}\varepsilon_{\rm AW}(c),c)
\]
simultaneously preserves the modular coordinate and the history. This component of the
construction is realized and does not use the leading term of $L(E,s)$. The
squares
\begin{equation}\label{eq:pdf2-bsd-pullback-base-change}
 P_{E,\mathbb Z}^{\rm gen}\otimes_{\mathbb Z}^{\mathbf L}
      \mathbb Z/\ell^n
 \xrightarrow{\ \sim\ }P_{E,\ell,n}^{\rm gen},
 \qquad
 \widehat\Pi_{0,\mathbb Z}\otimes1=\widehat\Pi_{0,\ell,n}
\end{equation}
they prove that unit, pullback, and counit are the same for all coefficients.

The two integral publications and their base changes
\begin{equation}\label{eq:pdf2-bsd-copub}
 P_{E,R}=(P_{E,R}^K,P_{E,R}^{\rm PT}),
 \qquad P_{E,R}^K(1_E)=z_{K,R},\quad
 P_{E,R}^{\rm PT}(1_E)=z_{{\rm PT},R}
\end{equation}
must follow from the same unit $1_E$. The reduced duality is typed by
\[
 D_3(M):=\operatorname{RHom}_\Lambda(M^\iota,\Lambda)[-3],
 \qquad
 X^{\rm orth}:C_E^{\rm red}\xrightarrow{\sim}D_3(C_E^{\rm red}).
\]
With $L_{{\rm root},\ell,n}=A_{\ell,n}z_{{\rm PT},\ell,n}$ and the complex projector
$p_{\rm root}$, the based normalization gives
\begin{equation}\label{eq:pdf2-bsd-raiz}
 p_{\rm root}z_{K,\ell,n}=z_{{\rm PT},\ell,n}.
\end{equation}

\section{Local filler and global decomposition without erasing the finite lattice}

The decomposition is first written over the integral owner. For
each finite--singular label $\lambda=(\ell,m,n,\ldots)$ there is an
incidence coefficient $a_\lambda=\ell c_\lambda\in\mathbb Z$ fixed by
$\operatorname{Rel}_{\rm det}$; in the corresponding fiber
$A_{\ell,n}$ acts by reduction of that same integer. The integral
free block with primitive $g$ is
\begin{equation}\label{eq:pdf2-bsd-bloque-normal}
 \partial f=a_\lambda e+b,\qquad
 \partial e=g,\qquad
 \partial b=-a_\lambda g.
\end{equation}
Tensoring it with the fiber indicated by $\lambda$ returns literally
$a_\lambda=\ell c_\lambda$. Consequently, if
$z_K=ar+ue+vb$ is a cycle and $z_{\rm PT}=r$, then
$\partial z_K=(u-a_\lambda v)g=0$ implies $u=a_\lambda v$. The root equality
\eqref{eq:pdf2-bsd-raiz} gives $a=1$ and, therefore,
\begin{equation}\label{eq:pdf2-bsd-filler-local}
 h_*:=-vf,\qquad \partial h_*=z_{\rm PT}-z_K.
\end{equation}

The multisection $\Sigma_{\rm det}^{\rm mult}$ labels each generator by
\begin{equation}\label{eq:pdf2-bsd-etiqueta-completa}
 \lambda=(\ell,m,n,\gamma,\operatorname{or},\operatorname{Carr},
 \operatorname{Mem},\partial,\tau^{\rm term},\operatorname{loc}).
\end{equation}
The based order separates three disjoint and exhaustive classes: the root unit, the finite--singular blocks with terminal $h_*$, and the cells of the local--global integral lattice. Since $\partial_{\rm det}$ preserves all labels except for the descent prescribed by carry, these classes define projectors of complexes $p_{\rm root}$, $p_{\rm FS}$, and $p_{\rm lat}$ with
\begin{equation}\label{eq:pdf2-bsd-tres-proyectores}
 p_{\rm root}+p_{\rm FS}+p_{\rm lat}=I,
 \qquad p_ip_j=\delta_{ij}p_i.
\end{equation}

\begin{theorem}[Global decomposition based]
\label{thm:pdf2-bsd-descomposicion-global}
There exists an isomorphism of complexes compatible with places, orientation,
corestriction, and duality
\begin{equation}\label{eq:pdf2-bsd-descomposicion-global}
 \Theta_E:\mathcal C_E^{\rm loc,red}
 \xrightarrow{\ \sim\ }
 \mathcal C_E^{\rm FS,ctr}\oplus\mathcal L_E,
\end{equation}
where
\begin{equation}\label{eq:pdf2-bsd-sumando-fs-y-lattice}
 \begin{aligned}
 \mathcal C_E^{\rm FS,ctr}
  &:=\bigoplus_{\lambda\in\Lambda_E^{\rm FS}}
   [\mathbb Zf_\lambda\to
    \mathbb Ze_\lambda\oplus\mathbb Zb_\lambda
                    \to\mathbb Zg_\lambda],\\
 \mathcal L_E&:=[M_E\xrightarrow{D_E}N_E].
 \end{aligned}
\end{equation}
Each block of the first summand has the integral differential
\eqref{eq:pdf2-bsd-bloque-normal}; its fiber $(\ell,n)$ is obtained through
$-\otimes_{\mathbb Z}^{\mathbf L}A_{\ell,n}$. The complex $\mathcal L_E$
preserves, without
contracting them, the Smith divisors and the
Tate--Shafarevich--Tamagawa components.
\end{theorem}

\begin{proof}
The generators are ordered by the label
\eqref{eq:pdf2-bsd-etiqueta-completa}. The common unit occupies the first
block and is removed by $q_{\rm root}$. For each terminal
$(z_{\rm PT},h_*)$, the relations of $\operatorname{Rel}_{\rm det}$
identify exactly four generators
$(f_\lambda,e_\lambda,b_\lambda,g_\lambda)$; the incidence matrix is
\eqref{eq:pdf2-bsd-bloque-normal}, with $g_\lambda$ primitive because its
arrival label occurs only once in the oriented basis. The remaining
cells have a local--global label and form the two free networks
$M_E,N_E$. There is no fourth class: the multisection is complete, and
$\operatorname{Surv}_{\rm det}$ requires every history to have either an FS terminal
or a network coordinate. The change of basis is unitriangular with
diagonal one, since a boundary introduces only earlier labels. This
constructs $\Theta_E$ and simultaneously proves its exhaustiveness and its
compatibility with the bases. The duality $X^{\rm orth}$ exchanges the
two networks and preserves each FS block, so it also intertwines
$\Theta_E$.
\end{proof}

On a normal block, define
\begin{equation}\label{eq:pdf2-bsd-hfs-generadores}
 H_\lambda(g_\lambda)=e_\lambda,
 \qquad H_\lambda(b_\lambda)=f_\lambda,
 \qquad H_\lambda(e_\lambda)=H_\lambda(f_\lambda)=0.
\end{equation}
A calculation on the four generators gives
$\partial H_\lambda+H_\lambda\partial=I$. Therefore
\begin{equation}\label{eq:pdf2-bsd-hfs-global}
 H_{\rm FS}:=\Theta_E^{-1}
 \left(\bigoplus_{\lambda\in\Lambda_E^{\rm FS}}H_\lambda\oplus0\right)
 \Theta_E,
 \qquad
 \partial H_{\rm FS}+H_{\rm FS}\partial=p_{\rm FS}.
\end{equation}
The terminal coordinate of $\mathfrak G_{\rm det}(E)$ ensures
$q_{\rm root}z_K\in\operatorname{im}p_{\rm FS}$; then
$h_*=-H_{\rm FS}(q_{\rm root}z_K)$ reproduces
\eqref{eq:pdf2-bsd-filler-local}. It is essential that the right-hand side of
\eqref{eq:pdf2-bsd-hfs-global} be $p_{\rm FS}$, not $q_{\rm root}$:
contracting the lattice $\mathcal L_E$ would erase the finite factors that BSD
must publish.

\section{Integral owner and base change to all primes}

The preceding decomposition is executed first in the integral basis of
histories, not in a primary fiber. For
\[
 R\in\{\mathbb Z,\mathbb Z_\ell,\mathbb Z/\ell^n,
          \mathbb Q_\ell\}
\]
are fixed
\begin{equation}\label{eq:pdf2-bsd-complejos-base-change}
 \mathcal C_{E,R}^{\rm loc,red}:=
 \mathcal C_{E,\mathbb Z}^{\rm loc,red}
       \otimes_{\mathbb Z}^{\mathbf L}R,
 \qquad
 \mathcal L_{E,R}:=
 [M_E\otimes R\xrightarrow{D_E\otimes1}N_E\otimes R].
\end{equation}
The unitriangular order of labels used in
\cref{thm:pdf2-bsd-descomposicion-global} has integer coefficients. Consequently,
\begin{equation}\label{eq:pdf2-bsd-theta-base-change}
 \Theta_{E,R}=\Theta_{E,\mathbb Z}\otimes1,
 \qquad
 H_{{\rm FS},R}=H_{{\rm FS},\mathbb Z}\otimes1,
 \qquad
 \partial H_{{\rm FS},R}+H_{{\rm FS},R}\partial=p_{{\rm FS},R}.
\end{equation}
In particular, the $3$-primary fibre is one of these specializations and cannot
conceal the fibres $\ell\ne3$.

Before forming a finite cokernel, the rational inverse is constructed. The ordering
of the terminals of $\Sigma_{\rm det}^{\rm mult}$ gives bases
$(m_1,\ldots,m_t)$ of $M_E$ and $(n_1,\ldots,n_t)$ of $N_E$, and a based bijection
$\sigma_E$ between the two lists. If
\[
 d_j:=\bigl\langle n_{\sigma_E(j)}^\vee,D_Em_j\bigr\rangle,
\]
the terminal relation and the duality $X^{\rm orth}$ yield $d_j\ne0$ before
dimensions are taken. Define
\begin{equation}\label{eq:pdf2-bsd-d0-nilpotente}
 \begin{aligned}
  D_{E,0}:M_E\otimes\mathbb Q&\longrightarrow N_E\otimes\mathbb Q,
  &D_{E,0}m_j&=d_jn_{\sigma_E(j)},\\
  N_E^{\rm tri}&:=D_{E,0}^{-1}(D_E\otimes\mathbb Q)-I.
 \end{aligned}
\end{equation}
A boundary contains only its terminal element and strictly preceding labels;
therefore $N_E^{\rm tri}$ is strictly triangular in the ordered list and
$(N_E^{\rm tri})^t=0$. The operator
\begin{equation}\label{eq:pdf2-bsd-j-e-inverso-racional}
 \boxed{
 J_E:=\left(\sum_{r=0}^{t-1}(-N_E^{\rm tri})^r\right)D_{E,0}^{-1}:
 N_E\otimes\mathbb Q\longrightarrow M_E\otimes\mathbb Q }
\end{equation}
is therefore computed using finitely many operations on the
incidence matrix. Since
$D_E\otimes\mathbb Q=D_{E,0}(I+N_E^{\rm tri})$, the finite geometric sum
proves both identities
\begin{equation}\label{eq:pdf2-bsd-j-e-dos-lados}
 J_E(D_E\otimes\mathbb Q)=I_{M_E\otimes\mathbb Q},
 \qquad
 (D_E\otimes\mathbb Q)J_E=I_{N_E\otimes\mathbb Q}.
\end{equation}
This is the rational contraction of the residual lattice; it uses neither rank,
$\operatorname{Sha}(E)$, group orders, nor the main term.

The integer matrix $D_E$ preserves all its elementary divisors. For each
prime, one defines, without inferring it from the $3$-primary fibre,
\begin{equation}\label{eq:pdf2-bsd-f-ell-def}
 D_{E,\ell}:=D_E\otimes\mathbb Z_\ell,
 \qquad
 F_{E,\ell}:=\operatorname{coker}D_{E,\ell}.
\end{equation}
The two identities of \eqref{eq:pdf2-bsd-j-e-dos-lados} imply that
$F_E:=\operatorname{coker}D_E$ is finite. Therefore its
primary decomposition is finite and exact:
\begin{equation}\label{eq:pdf2-bsd-f-descomposicion-primaria}
 F_E=\operatorname{coker}D_E
 \xrightarrow{\ \sim\ }
 \bigoplus_{\ell}F_{E,\ell},
 \qquad
 [n]\longmapsto([n\otimes1]_\ell)_\ell.
\end{equation}
The inverse is obtained generator by generator by the Chinese remainder theorem
applied to each Smith divisor. This is the edge that makes it possible to publish
$\operatorname{Sha}(E)$ and all Tamagawa factors simultaneously.

\section{Bocksteins on all pages}

Let
\begin{equation}\label{eq:pdf2-bsd-s-universal}
 S_{E,0}(T):V_0[[T]]\longrightarrow W_0[[T]]
\end{equation}
the rational based differential constructed by the integral matrices of the
determinantal reader. For each prime $\ell$ and for the complex chart, one
obtains by change of basis
\begin{equation}\label{eq:pdf2-bsd-s-todos-ell}
 S_{E,\ell}(T):=S_{E,0}(T)\otimes_{\mathbb Q}\mathbb Q_\ell,
 \qquad
 S_{E,\mathbb C}(T):=S_{E,0}(T)\otimes_{\mathbb Q}\mathbb C.
\end{equation}
It is invertible over $\mathbb Q((T))$ and, therefore, over each
$\mathbb Q_\ell((T))$ and $\mathbb C((T))$. We shall write $S_E(T)$ when the
basis is irrelevant. Since $K[[T]]$ is a discrete valuation ring for
$K=\mathbb Q,\mathbb Q_\ell$ or $\mathbb C$, its Smith form is
\begin{equation}\label{eq:pdf2-bsd-smith-t-adico}
 U(T)S_E(T)V(T)=\operatorname{diag}
 (T^{a_1},\ldots,T^{a_s},1,\ldots,1),
 \qquad a_i\ge1.
\end{equation}
The exponents $a_i$ are the same in every fiber: they are the valuations
in $T$ of the nonzero minors of the rational matrix
\eqref{eq:pdf2-bsd-s-universal}; replacing $\mathbb Q$ by
$\mathbb Q_\ell$ or $\mathbb C$ does not change those valuations.
One obtains
\begin{equation}\label{eq:pdf2-bsd-orden-smith}
 d:=\operatorname{ord}_T\det S_E(T)
 =\sum_{i=1}^sa_i
 =\sum_{q\ge1}\dim_K V_q.
\end{equation}
The filtration produces $A_q=V_q/V_{q+1}$,
$B_q=\operatorname{im}\beta_q$ and isomorphisms
\begin{equation}\label{eq:pdf2-bsd-bockstein-reducido}
 \bar\beta_q:A_q\xrightarrow{\sim}B_q,
 \qquad
 j_q:B_q\xrightarrow{\sim}A_q^\vee\otimes(I^q/I^{q+1}).
\end{equation}
With the regular page $\bar S_E(0)$ and the Koszul signs, all the
pages glue to
\begin{equation}\label{eq:pdf2-bsd-rboc-lt}
 R_{\rm Boc}^{\rm der}(S_E)
 =\tau_{\rm reg}\otimes\bigotimes_{q\ge1}\det(\bar\beta_q)
 =\left[T^{-d}\det S_E(T)\right]_{T=0}.
\end{equation}
No page is replaced by the total order $d$.

\section{Complete Kato--FERS table before the determinant}

For a nondegenerate history
$b_\lambda=[x_0\to\cdots\to x_i]$ let
$b_{\lambda,\le j}=[x_0\to\cdots\to x_j]$ and
$b_{\lambda,\ge j}=[x_j\to\cdots\to x_i]$. The complete weight of the
carry reader is
\begin{equation}\label{eq:pdf2-bsd-peso-t}
 \nu_T(\gamma):=\operatorname{dep}(\gamma)
                 +\operatorname{Carr}(\gamma),
\end{equation}
so that $\nu_T([x_i])=0$ and every normalized route of positive length
has $\nu_T\ge1$. The basis $\mathcal B_i^K$ contains one row for each
complete label $\lambda$; the basis $\mathcal B_i^{\rm Iw}$ contains the
generator $\upsilon_i(b_\lambda)$ with the same faces, sheet, orientation,
memory, and terminal. Therefore, $\upsilon_i$ is a bijection written from the
same list of labels, not chosen by means of a subsequent count.

The string map is
\begin{equation}\label{eq:pdf2-bsd-phi-kato}
 \Phi_{K/{\rm FERS}}:=
 \mu_{\rm Iw}\circ(P_E^K\widehat\otimes\operatorname{car}_T)
 \circ\Delta_{\rm AW}:
 \mathcal C_E^K\longrightarrow\mathcal C_E^{\rm Iw}.
\end{equation}
Its exhaustive table, parametrized by all $i$, all labels, and
all cuts, is
\begin{equation}\label{eq:pdf2-bsd-tabla-kato-fers}
\begin{array}{c|c|c|c}
 \text{row}&\text{cut}&\text{image}&\text{weight }T\\ \hline
 b_\lambda&j=i&\upsilon_i(b_\lambda)&0\\
 b_\lambda&0\le j<i&
 \epsilon_{\lambda,j}\,
 \mu_{\rm Iw}\!\left(P_E^K(b_{\lambda,\le j}),
       \operatorname{car}_T(b_{\lambda,\ge j})\right)&
 \nu_T(b_{\lambda,\ge j})\ge1\\
 \text{degenerate simplex}&\text{any }j&0&--
\end{array}
\end{equation}
where $\epsilon_{\lambda,j}$ is the sign already fixed by
$\operatorname{or}_{\rm det}$. Equivalently,
\begin{equation}\label{eq:pdf2-bsd-expansion-kato-fers}
 \Phi_{K/{\rm FERS}}^i(b_\lambda)
 =\upsilon_i(b_\lambda)+
  \sum_{j=0}^{i-1}T^{\nu_T(b_{\lambda,\ge j})}
  \epsilon_{\lambda,j}\,
  \mu_{\rm Iw}\!\left(P_E^K(b_{\lambda,\le j}),
          \widehat{\operatorname{car}}_T(b_{\lambda,\ge j})\right).
\end{equation}

\begin{theorem}[Kato--FERS-based equivalence]
\label{thm:pdf2-bsd-equivalencia-kato-fers}
The map \eqref{eq:pdf2-bsd-phi-kato} is a filtered equivalence of
perfect complexes, preserves the unit, and, in all complete bases,
\begin{equation}\label{eq:pdf2-bsd-i-plus-tb}
 [\Phi_{K/{\rm FERS}}^i](T)=I+TB_i(T),
 \qquad
 [\Phi_{K/{\rm FERS}}^i]^{-1}
 =\sum_{r\ge0}(-TB_i(T))^r.
\end{equation}
In particular,
\begin{equation}\label{eq:pdf2-bsd-rho-uno}
 \rho_E(T):=
 \prod_i\det(I+TB_i(T))^{(-1)^i},
 \qquad \rho_E(0)=1,
 \qquad
 \rho_{E,\ell}:=\rho_E\otimes_{\mathbb Q}\mathbb Q_\ell.
\end{equation}
\end{theorem}

\begin{proof}
The Alexander--Whitney simplicial identity, the compatibility of $P_E^K$ with faces, and the additivity of $\nu_T$ under concatenation prove that $\Phi_{K/{\rm FERS}}$ commutes with the differentials. In every row of \eqref{eq:pdf2-bsd-tabla-kato-fers}, the sole zero-weight term is the cut $j=i$; all the others belong to $T\Lambda$. Since $\upsilon_i$ pairs the complete list of histories with itself, reduction modulo $T$ is the identity on every generator. This yields \eqref{eq:pdf2-bsd-i-plus-tb}; completeness of $K[[T]]$ makes the inverse converge. For $i=0$, the table contains only the unit, so the root is normalized before a leading term is evaluated.
\end{proof}

The equivalence preserves separately the ledger
\begin{equation}\label{eq:pdf2-bsd-ledger-diez}
 \mathcal J_{\rm det}=\{\deg,\operatorname{Frob},\operatorname{Ner},
 K,\operatorname{Boc},\operatorname{Sha},\operatorname{Tam},
 \operatorname{tors}^{-2},\Omega,\operatorname{exc}\}.
\end{equation}
Those ten labels respectively record: the total weight; the Euler norm;
the Manin--Neron base change; the finite--singular relation; all the
pages $j_q\bar\beta_q$; the Kummer cone; the component quotients;
the torsional network; the Neron integral; and the exceptional orientation. The
table permits no cancellations between labels.

\section{Poitou--Tate Map from Mordell--Weil to the Full Flag}

First, the Mordell--Weil lattice is preserved and only afterwards is it rationalized:
\begin{equation}\label{eq:pdf2-bsd-mw-libre}
 M_{E,\mathbb Z}^{\rm MW}:=E(\mathbb Q)/E(\mathbb Q)_{\rm tors},
 \qquad
 M_E^{\rm MW}:=M_{E,\mathbb Z}^{\rm MW}\otimes_{\mathbb Z}K,
 \qquad K=\mathbb Q.
\end{equation}
$\operatorname{Surv}_{{\rm det},\mathbb Z}^{\rm MW}$ denotes the free submodule
generated by the integral terminal rows of the
Mordell--Weil summand.
The Kummer lift with complete history is
\begin{equation}\label{eq:pdf2-bsd-kummer-hmt}
 \begin{aligned}
 \kappa_{E,\mathbb Z}^{\rm HMT}:M_{E,\mathbb Z}^{\rm MW}
   &\longrightarrow
   Z^1(\mathcal C_{E,\mathbb Z}^{\rm red})
   \cap\operatorname{Surv}_{{\rm det},\mathbb Z}^{\rm MW},\\
 \kappa_E^{\rm HMT}&:=
   \kappa_{E,\mathbb Z}^{\rm HMT}\otimes_{\mathbb Z}\mathbb Q:
   M_E^{\rm MW}\longrightarrow Z^1(\mathcal C_E^{\rm red}).
 \end{aligned}
\end{equation}
and $\operatorname{pg}_q$ publishes the coordinate of page $V_q$. The
extension $\operatorname{Ext}_{\Gamma,{\rm det}}$, the multisection, and
survival construct, for each $q$, a based gluing
\begin{equation}\label{eq:pdf2-bsd-glue-q}
 \operatorname{gl}_q:V_q\longrightarrow
 \operatorname{Surv}_{\rm det}^{\rm MW}
\end{equation}
which preserves sheet, carry, memory, and terminal. Its equations of
unit--counit are
\begin{equation}\label{eq:pdf2-bsd-unidad-counidad-paginas}
 \operatorname{pg}_rP_E^{\rm PT}\operatorname{gl}_q
 =\delta_{rq}I_{V_q},
 \qquad
 \sum_{q\ge1}\operatorname{gl}_q
     \operatorname{pg}_qP_E^{\rm PT}=I
 \quad\text{in the surviving MW summand}.
\end{equation}
The first equality is checked label by label: the level
$q$ gluing has carry $q$, zero on the other pages, and the terminal of that same
history. The second is the exhaustiveness of
$\Sigma_{\rm det}^{\rm mult}$; the sum is finite because the Bockstein tower
terminates.

For this submodule, set it to
\begin{equation}\label{eq:pdf2-bsd-redes-paginas-saturadas}
 V_{q,\mathbb Z}:=\operatorname{pg}_qP_E^{\rm PT}
   (\operatorname{Surv}_{{\rm det},\mathbb Z}^{\rm MW}),
 \qquad
 A_{q,\mathbb Z}:=V_{q,\mathbb Z}/V_{q+1,\mathbb Z}.
\end{equation}
In the terminal bases, $\operatorname{pg}_q$ and
$\operatorname{gl}_q$ have entries $0,\pm1$, and their two compositions are
those of \eqref{eq:pdf2-bsd-unidad-counidad-paginas}. Therefore,
$V_{q+1,\mathbb Z}$ is a based direct summand of $V_{q,\mathbb Z}$,
$A_{q,\mathbb Z}$ is free and saturated, and
$V_q=V_{q,\mathbb Z}\otimes\mathbb Q$,
$A_q=A_{q,\mathbb Z}\otimes\mathbb Q$.
Let $s_q:A_{q,\mathbb Z}\to V_{q,\mathbb Z}$ be the section of the
terminal basis. Then
$V_{k,\mathbb Z}=\bigoplus_{q\geq k}s_q(A_{q,\mathbb Z})$. The splitting
of jets that preserves the $q$ occurrences of a depth-$q$ row is
\begin{equation}\label{eq:pdf2-bsd-desdoblamiento-jets}
 \begin{aligned}
 \operatorname{Jet}_{\rm Boc,\mathbb Z}(E)
   &:=\bigoplus_{q\geq1}\bigoplus_{k=1}^{q}
       A_{q,\mathbb Z}\,e_{q,k},\\
 \mathcal U_{\rm jet}:
 \operatorname{Jet}_{\rm Boc,\mathbb Z}(E)
   &\xrightarrow{\ \sim\ }
     \bigoplus_{k\geq1}V_{k,\mathbb Z},\\
 \mathcal U_{\rm jet}((a_{q,k})_{q,k})
   &:=(v_k)_k,
   \qquad v_k:=\sum_{q\geq k}s_q(a_{q,k}).
 \end{aligned}
\end{equation}
The inverse decomposes each $v_k$ into the preceding direct sum. In particular,
\begin{equation}\label{eq:pdf2-bsd-dimension-jet}
 \operatorname{rank}_{\mathbb Z}\operatorname{Jet}_{\rm Boc,\mathbb Z}(E)
 =\sum_q q\,\operatorname{rank}_{\mathbb Z}A_{q,\mathbb Z}
 =\sum_k\operatorname{rank}_{\mathbb Z}V_{k,\mathbb Z}=d.
\end{equation}
The gluing restricts, in those same bases, to
$\operatorname{gl}_{q,\mathbb Z}:V_{q,\mathbb Z}\to
\operatorname{Surv}_{{\rm det},\mathbb Z}^{\rm MW}$.

They are defined, without using dimensions,
\begin{equation}\label{eq:pdf2-bsd-psi-y-lambda-q}
 \begin{aligned}
  \psi_{q,\mathbb Z}&:=
      \operatorname{pg}_qP_E^{\rm PT}\kappa_{E,\mathbb Z}^{\rm HMT}:
      M_{E,\mathbb Z}^{\rm MW}\to V_{q,\mathbb Z},\\
  \lambda_{q,\mathbb Z}&:=
      \operatorname{pr}_{\rm MW}P_E^{\rm Man}
      \operatorname{gl}_{q,\mathbb Z}:
      V_{q,\mathbb Z}\to M_{E,\mathbb Z}^{\rm MW},\\
  \psi_q&:=\operatorname{pg}_qP_E^{\rm PT}\kappa_E^{\rm HMT}:
      M_E^{\rm MW}\to V_q,\\
  \lambda_q&:=\operatorname{pr}_{\rm MW}P_E^{\rm Man}
      \operatorname{gl}_q:V_q\to M_E^{\rm MW}.
 \end{aligned}
\end{equation}
The intermediate domains remain linked by the two based compatibilities
that form part of the Poitou--Tate gluing and that are written here over their
reachable action:
\begin{equation}\label{eq:pdf2-bsd-kummer-glue-compatibilidad}
 \kappa_E^{\rm HMT}\operatorname{pr}_{\rm MW}P_E^{\rm Man}
       \operatorname{gl}_q=\operatorname{gl}_q,
 \qquad
 \operatorname{pr}_{\rm MW}P_E^{\rm Man}\kappa_E^{\rm HMT}
       =I_{M_E^{\rm MW}}.
\end{equation}
The same two identities, with subscript $\mathbb Z$, hold over
$M_{E,\mathbb Z}^{\rm MW}$ and $V_{q,\mathbb Z}$, because all the preceding matrices
are integral and based.

\begin{theorem}[Poitou--Tate Flag-Based Isomorphism]
\label{thm:pdf2-bsd-psi-pt}
The maps
\begin{equation}\label{eq:pdf2-bsd-psi-pt}
 \begin{aligned}
 \Psi_{\rm PT}:M_E^{\rm MW}&\longrightarrow
       \operatorname{Flag}_{\rm Boc}(E):=\bigoplus_{q\ge1}V_q,
 &P&\longmapsto(\psi_q(P))_{q\ge1},\\
 \Psi_{\rm PT}^{-1}:\operatorname{Flag}_{\rm Boc}(E)&\longrightarrow M_E^{\rm MW},
 &(v_q)_q&\longmapsto\sum_{q\ge1}\lambda_q(v_q)
 \end{aligned}
\end{equation}
are inverses. Consequently,
\begin{equation}\label{eq:pdf2-bsd-psi-pt-integral}
 \Psi_{{\rm PT},\mathbb Z}:M_{E,\mathbb Z}^{\rm MW}
 \xrightarrow{\ \sim\ }\bigoplus_{q\ge1}V_{q,\mathbb Z},
 \qquad
 \Psi_{{\rm PT},\mathbb Z}^{-1}((v_q)_q)=
 \sum_q\lambda_{q,\mathbb Z}(v_q),
\end{equation}
and, only after extending the scale,
\begin{equation}\label{eq:pdf2-bsd-rango-orden}
 \operatorname{rank}E(\mathbb Q)
 =\sum_{q\ge1}\dim_K V_q
 =d=\operatorname{ord}_T\det S_E(T).
\end{equation}
\end{theorem}

\begin{proof}
Applying \eqref{eq:pdf2-bsd-unidad-counidad-paginas} together with \eqref{eq:pdf2-bsd-kummer-glue-compatibilidad} yields, with all intermediate domains visible, $\psi_r\lambda_q=\delta_{rq}I$ and $\sum_q\lambda_q\psi_q=I_{M_E^{\rm MW}}$. The two formulas for \eqref{eq:pdf2-bsd-psi-pt} are therefore inverse before dimensions are compared. The terminal tables are integral and based; the same computation, without denominators, proves \eqref{eq:pdf2-bsd-psi-pt-integral}. Thus no sublattice index appears when a Mordell--Weil basis is subsequently chosen. Only then are dimensions taken and \eqref{eq:pdf2-bsd-orden-smith} used. The pairing $j_q\bar\beta_q$ transports the saturated bases to the height rows whose comparison with Neron--Tate is written in \eqref{eq:pdf2-bsd-altura-nt-generadores}.
\end{proof}

\section{Localization, Smith Range and Sha--Tamagawa Accuracy}

In this section, $v$ ranges over the finite places; at those of good reduction,
we have $\Phi_v(\mathbb F_v)=0$ and $c_v=1$, so that the local sums and products
have finite support. The infinite place is reserved for evaluating
the Neron differential in the fourth arrow. For each finite place $v$,
we have the triangle
\[
 \mathcal C_{E,v}^{\rm fin}\to\mathcal C_{E,v}\to
 \mathcal C_{E,v}^{\rm sing}\to\mathcal C_{E,v}^{\rm fin}[1]
\]
and the string map
\begin{equation}\label{eq:pdf2-bsd-ell-e}
 \ell_E(x,(y_v)_v)=(\operatorname{res}_vx-i_vy_v)_v.
\end{equation}
The reduced cone is
\begin{equation}\label{eq:pdf2-bsd-cono-reducido}
 \mathcal C_E^{\rm loc,red}:=
 q_{\rm root}\operatorname{Cone}(\ell_E)[-1].
\end{equation}
In cocycle degrees, an element is written $(x,(y_v),(h_v))$ and
\begin{equation}\label{eq:pdf2-bsd-diferencial-cono}
 d(x,(y_v),(h_v))=
 (dx,(dy_v),(\operatorname{res}_vx-i_vy_v-dh_v)_v).
\end{equation}

After separating by \eqref{eq:pdf2-bsd-psi-pt} the free Mordell--Weil lattice
and its dual, the FS blocks are contracted by
\eqref{eq:pdf2-bsd-hfs-global}. In the remaining lattice, the inverse
\eqref{eq:pdf2-bsd-j-e-inverso-racional} and its two checks
\eqref{eq:pdf2-bsd-j-e-dos-lados} give, without inferring acyclicity solely from
duality,
\begin{equation}\label{eq:pdf2-bsd-rango-completo}
 D_E\otimes\mathbb Q:M_E\otimes\mathbb Q
 \xrightarrow{\sim}N_E\otimes\mathbb Q.
\end{equation}
The integral normal form gives
\begin{equation}\label{eq:pdf2-bsd-smith-finito}
 UD_EV=\operatorname{diag}(s_1,\ldots,s_t),
 \qquad
 F_E:=\operatorname{coker}D_E,
 \qquad |F_E|=\prod_{i=1}^t|s_i|<\infty.
\end{equation}

Let
\begin{equation}\label{eq:pdf2-bsd-pi-componentes}
 \pi_\Phi:N_E\longrightarrow\bigoplus_v\Phi_v(\mathbb F_v)
\end{equation}
the reduction of each local coordinate modulo
$\mathcal C_{E,v}^{\rm fin}+d\mathcal C_{E,v}$; one has
$\pi_\Phi D_E=0$. For a class
$[\xi]\in\operatorname{Sha}(E)$, choose its finite primitives
$y_v$ and homotopies $h_v$ with
$\operatorname{res}_v\xi-i_vy_v=dh_v$. The formula
\begin{equation}\label{eq:pdf2-bsd-iota-sha-cociclos}
 \begin{aligned}
 n_\xi&:=\operatorname{pr}_{N_E}\operatorname{pr}_{\mathcal L_E}
       \Theta_E(\xi,(y_v),(h_v)),\\
 \iota_{\rm Sha}[\xi]&:=[n_\xi]\in F_E
 \end{aligned}
\end{equation}
is independent of the choices by
\eqref{eq:pdf2-bsd-diferencial-cono}.

To render the final surjectivity material, let
$\widetilde\phi_v$ be the based terminal representative of
$\phi_v\in\Phi_v(\mathbb F_v)$. The sum of counits is corrected at the
common unit and lifted by the structural operator:
\begin{equation}\label{eq:pdf2-bsd-lift-componentes}
 \operatorname{Lift}_\Phi((\phi_v)_v):=
 \operatorname{pr}_{N_E}\operatorname{Ext}_{\Gamma,{\rm det}}
 \left((\widetilde\phi_v)_v,
 -\widehat\Pi_0\!\left(\sum_v
       \varepsilon_{\rm AW}(\widetilde\phi_v)\right)\right).
\end{equation}
The corestriction relation of
$\operatorname{Rel}_{\rm det}$ makes this expression a cocycle and gives
\begin{equation}\label{eq:pdf2-bsd-lift-dos-identidades}
 \pi_\Phi\operatorname{Lift}_\Phi=I,
 \qquad
 \operatorname{im}D_E=
 \ker\pi_\Phi\cap\ker(\text{global class}).
\end{equation}

\begin{theorem}[Finite Sha--Tamagawa exact sequence]
\label{thm:pdf2-bsd-sha-tamagawa}
The preceding cocycle maps induce
\begin{equation}\label{eq:pdf2-bsd-sha-tamagawa}
 0\longrightarrow\operatorname{Sha}(E)
 \xrightarrow{\ \iota_{\rm Sha}\ }F_E
 \xrightarrow{\ \partial_{\rm loc}\ }
 \bigoplus_v\Phi_v(\mathbb F_v)
 \longrightarrow0,
 \qquad
 \partial_{\rm loc}[n]:=\pi_\Phi(n).
\end{equation}
In particular, $\operatorname{Sha}(E)$ is finite and
\begin{equation}\label{eq:pdf2-bsd-cardinal-sts}
 |F_E|=|\operatorname{Sha}(E)|\prod_vc_v,
 \qquad c_v:=|\Phi_v(\mathbb F_v)|.
\end{equation}
\end{theorem}

\begin{proof}
The equality $\pi_\Phi D_E=0$ makes $\partial_{\rm loc}$ well defined. If \eqref{eq:pdf2-bsd-iota-sha-cociclos} is a boundary, its global component is a boundary; thus $\iota_{\rm Sha}$ is injective. A class of $F_E$ lies in $\ker\partial_{\rm loc}$ if and only if all its singular coordinates admit finite representatives; by \eqref{eq:pdf2-bsd-diferencial-cono}, this is equivalent to lying in the image of $\iota_{\rm Sha}$. Finally, \eqref{eq:pdf2-bsd-lift-dos-identidades} gives a preimage of every tuple of components and proves surjectivity. The adjacent free terms of the long exact sequence are precisely the two summands separated by $\Psi_{\rm PT}$ and $X^{\rm orth}$; the adjacent FS terms are contractible by $H_{\rm FS}$. Their vanishing has not been assumed. The finiteness of $F_E$ follows from \eqref{eq:pdf2-bsd-rango-completo}; injectivity gives the finiteness of $\operatorname{Sha}(E)$, and taking orders gives \eqref{eq:pdf2-bsd-cardinal-sts}.
\end{proof}

\begin{corollary}[Complete primary discharge and integral recomposition]
\label{cor:pdf2-bsd-sha-tamagawa-todos-ell}
For each prime $\ell$, base change of the cocycle maps yields
the exact sequence
\begin{equation}\label{eq:pdf2-bsd-sha-tamagawa-ell}
 0\longrightarrow\operatorname{Sha}(E)[\ell^\infty]
 \xrightarrow{\ \iota_{{\rm Sha},\ell}\ }F_{E,\ell}
 \xrightarrow{\ \partial_{{\rm loc},\ell}\ }
 \bigoplus_v\Phi_v(\mathbb F_v)[\ell^\infty]
 \longrightarrow0,
\end{equation}
where
\begin{equation}\label{eq:pdf2-bsd-mapas-ell-generadores}
 \iota_{{\rm Sha},\ell}([\xi]\otimes1)
   =[n_\xi\otimes1],
 \qquad
 \partial_{{\rm loc},\ell}([n\otimes1])
   =(\pi_{\Phi,v}(n)\otimes1)_v.
\end{equation}
Moreover, the sum of all the sequences
\eqref{eq:pdf2-bsd-sha-tamagawa-ell} is exactly
\eqref{eq:pdf2-bsd-sha-tamagawa}; it is not merely its $3$-primary part.
\end{corollary}

\begin{proof}
The matrix $D_E$, the cocycle $n_\xi$, the reduction $\pi_\Phi$, and
$\operatorname{Lift}_\Phi$ are integral by
\eqref{eq:pdf2-bsd-complejos-base-change}--
\eqref{eq:pdf2-bsd-theta-base-change}. Since $\mathbb Z_\ell$ is flat over
$\mathbb Z$, tensoring the finite exact sequence
\eqref{eq:pdf2-bsd-sha-tamagawa} preserves exactness. For a finite group
$A$ one has canonically
$A\otimes\mathbb Z_\ell=A[\ell^\infty]$; this identifies the three terms
and gives the formulas on generators
\eqref{eq:pdf2-bsd-mapas-ell-generadores}. The Smith form
\eqref{eq:pdf2-bsd-smith-finito} writes
$s_i=\prod_\ell\ell^{v_\ell(s_i)}$, so that
\[
 F_E\cong\bigoplus_\ell F_{E,\ell},\qquad
 \operatorname{Sha}(E)\cong
   \bigoplus_\ell\operatorname{Sha}(E)[\ell^\infty],\qquad
 \Phi_v(\mathbb F_v)\cong
   \bigoplus_\ell\Phi_v(\mathbb F_v)[\ell^\infty].
\]
Only finitely many primes and finitely many places intervene because all the groups
are quotients of the finite group $F_E$. Summing the primary sequences recovers
the integral one. Taking orders first prime by prime and then multiplying,
\begin{equation}\label{eq:pdf2-bsd-cardinal-sts-todos-ell}
 \prod_\ell|F_{E,\ell}|
 =\prod_\ell|\operatorname{Sha}(E)[\ell^\infty]|
   \prod_v\prod_\ell|\Phi_v(\mathbb F_v)[\ell^\infty]|
 =|\operatorname{Sha}(E)|\prod_vc_v.
\end{equation}
Thus all the coefficients and all the Tamagawa factors are materially discharged.
\end{proof}

\section{Based analytic--arithmetic comparator}

The analytic coordinate is not postulated from the sought order. Let
$\lambda_3=\log4$, corresponding to the based cyclotomic generator
$\gamma_3=4$, and define, on a disk $U_E$ centered at $1$,
\begin{equation}\label{eq:pdf2-bsd-t-e-construido}
 T_E(s):=\frac{\exp(\lambda_3(s-1))-1}{\lambda_3}
 =(s-1)+\frac{\lambda_3}{2}(s-1)^2+O((s-1)^3).
\end{equation}
By direct calculation,
\begin{equation}\label{eq:pdf2-bsd-t-e-normalizacion-probada}
 T_E(1)=0,
 \qquad T_E'(s)=\exp(\lambda_3(s-1)),
 \qquad T_E'(1)=1.
\end{equation}
The inverse function theorem provides, after reducing $U_E$ if
necessary, a biholomorphic chart $s\leftrightarrow T_E(s)$.
To compare all fibers, take
$\gamma_\ell=1+\ell$ if $\ell$ is odd and $\gamma_2=5$, and set
\begin{equation}\label{eq:pdf2-bsd-t-ell-construido}
 T_{E,\ell}(s):=
 \frac{\exp(\log(\gamma_\ell)(s-1))-1}{\log(\gamma_\ell)}.
\end{equation}
The charts are related by a unique formal automorphism
$T_{E,\ell}=\vartheta_{\ell,3}(T_E)$ with
$\vartheta_{\ell,3}(0)=0$ and $\vartheta_{\ell,3}'(0)=1$. Therefore, neither the
order at $T$ nor the leading coefficient changes upon passing from $3$ to $\ell$.

The analytic continuation that reaches this chart is not declared either. By the
modularity theorem, let
\[
 f_E(z)=\sum_{m\geq1}a_m(E)e^{2\pi imz}
\]
the normalized newform of weight $2$ and level $N_E$ associated with $E$, and let
$\varepsilon_E\in\{\pm1\}$ be the sign fixed by
\begin{equation}\label{eq:pdf2-bsd-transformacion-modular}
 g_E(y):=f_E\!\left(\frac{iy}{\sqrt{N_E}}\right),
 \qquad
 g_E(1/y)=\varepsilon_Ey^2g_E(y).
\end{equation}
For $\Re s>3/2$, termwise integration gives
\begin{equation}\label{eq:pdf2-bsd-mellin-semiplano}
 \Lambda_E(s):=N_E^{s/2}(2\pi)^{-s}\Gamma(s)L(E,s)
 =\int_0^\infty g_E(y)y^s\,\frac{dy}{y}.
\end{equation}
Split the integral into $(0,1)$ and $(1,\infty)$ and, in the first part,
make the change $y=1/t$. The second equality in
\eqref{eq:pdf2-bsd-transformacion-modular} yields the material formula
\begin{equation}\label{eq:pdf2-bsd-mellin-continuada}
 \Lambda_E^{\rm cont}(s)=
 \int_1^\infty g_E(y)
       \bigl(y^s+\varepsilon_Ey^{2-s}\bigr)\,\frac{dy}{y}.
\end{equation}
The function $g_E(y)$ decreases exponentially when $y\to\infty$; the integral of \eqref{eq:pdf2-bsd-mellin-continuada}, and all its derivatives in $s$, therefore converge uniformly on compact sets. It thus defines an entire function and coincides with \eqref{eq:pdf2-bsd-mellin-semiplano} where the Dirichlet series converges. Since $N_E^{s/2}(2\pi)^{-s}\Gamma(s)$ is holomorphic and nonzero near $s=1$, division by that factor constructs, without using its order of vanishing, the holomorphic germ of $L(E,s)$ in a disk $U_E$ centered at $1$.

To avoid mixing a $\ell$-adic fibre with a complex matrix, one first takes
the integer polynomial, independent of $\ell$,
\begin{equation}\label{eq:pdf2-bsd-polinomio-local-entero}
 P_{E,p}(X):=\det(I-X\operatorname{Frob}_p\mid V_\ell E^{I_p})
 \in\mathbb Z[X].
\end{equation}
For $p\nmid N_E$, let $\alpha_p,\beta_p$ be the roots of $P_{E,p}$. Take
$W_p=\mathbb C^2$ with orthonormal basis and the normal realization
\begin{equation}\label{eq:pdf2-bsd-realizacion-normal-frobenius}
 F_p^{\mathbb C}:=\operatorname{diag}(\alpha_p,\beta_p),
 \qquad |\alpha_p|=|\beta_p|=p^{1/2},
 \qquad \|F_p^{\mathbb C}\|_1=2p^{1/2}.
\end{equation}
For the bad primes one likewise chooses a complex realization of the
local polynomial; they are finite and are preserved as separate blocks. The
companion matrix is not used, since its norm would not be uniformly controlled by the
moduli of its eigenvalues.

Set $\mathscr H_E:=\widehat\bigoplus_{p\nmid N_E}W_p$ and, for $X>0$,
\begin{equation}\label{eq:pdf2-bsd-truncados-fredholm}
 K_{E,X}(s):=\bigoplus_{\substack{p\leq X\\p\nmid N_E}}
          p^{-s}F_p^{\mathbb C},
 \qquad
 L_X(E,s):=\det(I-K_{E,X}(s))^{-1}
       \prod_{p\mid N_E}P_{E,p}(p^{-s})^{-1}.
\end{equation}
If $K_0$ is a compact of $\{\Re s>3/2\}$ and
$\sigma_0:=\inf_{s\in K_0}\Re s$, then
\begin{equation}\label{eq:pdf2-bsd-cota-traza}
 \sup_{s\in K_0}\|K_{E,X'}(s)-K_{E,X}(s)\|_1
 \leq C_E\sum_{X<p\leq X'}p^{1/2-\sigma_0}
 \xrightarrow[X,X'\to\infty]{}0.
\end{equation}
Thus $K_{E,X}\to K_E$ in trace norm, locally uniformly, and the
continuity of the Fredholm determinant proves, rather than defines,
\begin{equation}\label{eq:pdf2-bsd-fredholm-euler}
 \det_F(I-K_E(s))^{-1}
 \prod_{p\mid N_E}P_{E,p}(p^{-s})^{-1}
 =\lim_{X\to\infty}L_X(E,s)=L(E,s).
\end{equation}

To pass from the infinite product to a finite complex, one first preserves the
operator, not merely its determinant. Let $\mathscr H_E^+$ and
$\mathscr H_E^-$ be two based copies of
$\mathscr H_E\oplus\bigoplus_{p\mid N_E}W_p$. In the half-plane
$\Re s>3/2$, define
\begin{equation}\label{eq:pdf2-bsd-operadores-euler-inversos}
 \begin{aligned}
  \mathscr A_{E,X}(s)&:=
   (I-K_{E,X}(s))^{-1}\oplus
   \bigoplus_{p\mid N_E}(I-p^{-s}F_p^{\mathbb C})^{-1},\\
  \mathscr A_E(s)&:=\lim_{X\to\infty}\mathscr A_{E,X}(s):
       \mathscr H_E^+\longrightarrow\mathscr H_E^-.
 \end{aligned}
\end{equation}
The resolvent identity and \eqref{eq:pdf2-bsd-cota-traza} give, for each
compact $K_0$ of the half-plane,
\begin{equation}\label{eq:pdf2-bsd-resolvente-convergencia-traza}
 \sup_{s\in K_0}
 \|\mathscr A_{E,X'}(s)-\mathscr A_{E,X}(s)\|_1\longrightarrow0.
\end{equation}
Let $\mathcal U_{\rm bas}:\mathscr H_E^+\to\mathscr H_E^-$ be the unitary identification that
sends each vector of the positive basis to the homonymous vector of the
negative basis. The following determinants mean those of
$\mathcal U_{\rm bas}^{-1}\mathscr A_{E,X}=I+\mathcal L^1$ and
$\mathcal U_{\rm bas}^{-1}\mathscr A_E=I+\mathcal L^1$. Then
\begin{equation}\label{eq:pdf2-bsd-det-a-x-l-x}
 \det_F(\mathcal U_{\rm bas}^{-1}\mathscr A_{E,X}(s))=L_X(E,s),
 \qquad
 \det_F(\mathcal U_{\rm bas}^{-1}\mathscr A_E(s))=L(E,s).
\end{equation}
The second equality is precisely the limit of
\eqref{eq:pdf2-bsd-fredholm-euler}; the sign is correct because the
inverse of $I-K_E$ is taken.

The finite corestriction is constructed from the terminals, without publishing
infinitely many primes in a finite matrix. Let $\mathcal I_E^+$ and
$\mathcal I_E^-$ be the terminal lists of the integral bases of $M_E$ and
$N_E$, respectively; they are finite by
\eqref{eq:pdf2-bsd-sumando-fs-y-lattice}. With $\sigma_*>3/2$ fixed, the
local--global intertwiner embeds them into the Euler space by
\begin{equation}\label{eq:pdf2-bsd-embedding-terminal-euler}
 \iota_{\rm Eul}^\pm(e_\eta):=
 \left(p^{-\sigma_*}\operatorname{res}_p^\pm
       \operatorname{Surv}_E(e_\eta)\right)_p
 \oplus\left(\operatorname{res}_p^\pm
       \operatorname{Surv}_E(e_\eta)\right)_{p\mid N_E}.
\end{equation}
The good sum is square--summable by the same bound as
\eqref{eq:pdf2-bsd-cota-traza}. Set
$\tau_\eta^\pm:=\iota_{\rm Eul}^\pm(e_\eta)$; independence of the terminal rows
gives the dual counits $\{\epsilon_{\eta}^{\pm}\}$, so that
$\epsilon_\eta^\pm(\tau_{\eta'}^\pm)=\delta_{\eta\eta'}$. The sums
\begin{equation}\label{eq:pdf2-bsd-proyecciones-terminales-finitas}
 P_E^\pm:=\sum_{\eta\in\mathcal I_E^\pm}
      \tau_\eta^\pm\epsilon_\eta^\pm,
 \qquad Q_E^\pm:=I-P_E^\pm
\end{equation}
are based finite-rank projections; their indices are fixed by the
terminal labels, not by the Mordell--Weil rank or by an order of
vanishing.

The coordinate $\deg$ of the ledger decomposes
\begin{equation}\label{eq:pdf2-bsd-graduacion-fredholm}
 \mathscr H_E^\pm=\widehat\bigoplus_i\mathscr H_E^{\pm,i},
 \qquad
 \mathcal I_E^{\pm,i}:=\{\eta\in\mathcal I_E^\pm:\deg\eta=i\},
 \qquad
 P_E^{\pm,i}:=
 \sum_{\eta\in\mathcal I_E^{\pm,i}}\tau_\eta^\pm\epsilon_\eta^\pm.
\end{equation}
The local relations preserve $\deg$ except for the differential $i\mapsto i+1$;
therefore
$\mathscr A_E=\bigoplus_i\mathscr A_E^i$ with
$\mathscr A_E^i:\mathscr H_E^{+,i}\to\mathscr H_E^{-,i+1}$. The four
blocks of \eqref{eq:pdf2-bsd-bloques-fredholm}, its Schur complement, and its truncations
are taken degree by degree, and by definition
\begin{equation}\label{eq:pdf2-bsd-schur-graduado}
 \mathscr D_E=\bigoplus_i\mathscr D_E^i,
 \qquad
 \mathscr D_E^i:=A_{00}^i-A_{01}^i(A_{11}^i)^{-1}A_{10}^i.
\end{equation}

The Fredholm continuation is written before it is used. In the Euler chart
$U_0$ set $\mathscr A_{E,0}:=\mathscr A_E$. The charts preserve a
common based complement $\mathscr Q_E^\pm$ and change only the finite family
of terminals. If
$\tau_{\eta,\alpha}^\pm:=
\operatorname{Ext}_{\Gamma,{\rm det},\alpha}(\tau_\eta^\pm)$ and
$\epsilon_{\eta,\alpha}^\pm$ is its dual basis, define
\begin{equation}\label{eq:pdf2-bsd-cambios-fredholm-cartas}
 C_{\alpha\beta}^\pm:=
 I_{\mathscr Q_E^\pm}\oplus
 \sum_{\eta\in\mathcal I_E^\pm}
 \tau_{\eta,\beta}^\pm\epsilon_{\eta,\alpha}^\pm,
 \qquad C_{\alpha\beta}^\pm-I\in\mathcal L^1,
 \qquad
 C_{\beta\gamma}^\pm C_{\alpha\beta}^\pm=C_{\alpha\gamma}^\pm.
\end{equation}
The difference from the identity has rank at most $|\mathcal I_E^\pm|$ and is
therefore trace class; holomorphy follows because each of the finitely many
terminal coordinates of $\operatorname{Ext}_{\Gamma,{\rm det},\alpha}$ is holomorphic. The final equality is
verified in $\mathscr Q_E^\pm$ and in every $\tau_{\eta,\alpha}^\pm$. It is defined
recursively on the overlaps
\begin{equation}\label{eq:pdf2-bsd-continuacion-fredholm-cartas}
 \mathscr A_{E,\beta}:=
 C_{\alpha\beta}^-\mathscr A_{E,\alpha}
 (C_{\alpha\beta}^+)^{-1}.
\end{equation}
The operators $C_{\alpha\beta}^\pm$ are holomorphic because every coefficient
is a coordinate of $\operatorname{Ext}_{\Gamma,{\rm det}}$; the cocycle of
\eqref{eq:pdf2-bsd-cambios-fredholm-cartas} proves that
\eqref{eq:pdf2-bsd-continuacion-fredholm-cartas} does not depend on the chain of
charts. This is the determinant--class family henceforth denoted by
$\mathscr A_E(s)$. With
$C_{0\alpha}^\pm:=C_{\alpha-1,\alpha}^\pm\cdots C_{0,1}^\pm$, one additionally transports
$P_{E,\alpha}^{\pm,i}:=
C_{0\alpha}^\pm P_E^{\pm,i}(C_{0\alpha}^\pm)^{-1}$ and
$\mathscr H_{E,\alpha}^{\pm,i}:=
C_{0\alpha}^\pm\mathscr H_E^{\pm,i}$; thus all Schur
$\mathscr D_{E,\alpha}^i$ are those of
\eqref{eq:pdf2-bsd-schur-graduado}, not new symbols. After reducing each
chart, the block
\begin{equation}\label{eq:pdf2-bsd-bloques-fredholm}
 A_{ab}(s):=P_a^-\mathscr A_E(s)P_b^+,
 \qquad P_0^\pm:=P_E^\pm,\quad P_1^\pm:=Q_E^\pm,
 \qquad A_{11}(s):Q_E^+\mathscr H_E^+\xrightarrow{\sim}
                       Q_E^-\mathscr H_E^-
\end{equation}
is invertible. This is the standard finite reduction of the Fredholm operator;
the specific choice of $P_E^\pm$ is given by
\eqref{eq:pdf2-bsd-proyecciones-terminales-finitas}. The basis also fixes
the unitary identification
$U_{11}:Q_E^+\mathscr H_E^+\to Q_E^-\mathscr H_E^-$; by construction,
$U_{11}^{-1}A_{11}=I+\mathcal L^1$, and only this operator is acted on by
$\det_F$. The copies $+$ and $-$ already constitute the even--odd totalization; no
second alternating sign is applied to the determinant of the inverse operator.

For each truncation that contains the support of those terminals, it is defined materially.
\begin{equation}\label{eq:pdf2-bsd-cor-lift-truncados}
 \begin{aligned}
  \operatorname{Cor}_{E,X}&:=P_E^-,\\
  \operatorname{Lift}_{E,X}(v)&:=
       v-A_{11,X}^{-1}A_{10,X}v,\\
  \mathscr D_{E,X}(s)&:=
       \operatorname{Cor}_{E,X}\mathscr A_{E,X}(s)
       \operatorname{Lift}_{E,X}
       =A_{00,X}-A_{01,X}A_{11,X}^{-1}A_{10,X}.
 \end{aligned}
\end{equation}
By block multiplication,
\begin{equation}\label{eq:pdf2-bsd-eliminacion-schur}
 \begin{pmatrix}I&-A_{01,X}A_{11,X}^{-1}\\0&I\end{pmatrix}
 \mathscr A_{E,X}
 \begin{pmatrix}I&0\\-A_{11,X}^{-1}A_{10,X}&I\end{pmatrix}
 =\begin{pmatrix}\mathscr D_{E,X}&0\\0&A_{11,X}\end{pmatrix}.
\end{equation}
Thus corestriction is substantive: its lift, its inverse on the complement,
and its finite differential are written explicitly. Convergence of the resolvent
implies
\begin{equation}\label{eq:pdf2-bsd-corestriccion-limite}
 \mathscr D_E(s):=\lim_{X\to\infty}\mathscr D_{E,X}(s)
 =A_{00}-A_{01}A_{11}^{-1}A_{10},
 \qquad
 \|\mathscr D_{E,X}-\mathscr D_E\|_1\to0.
\end{equation}
The determinant of the complement is not discarded: by
\eqref{eq:pdf2-bsd-eliminacion-schur},
\begin{equation}\label{eq:pdf2-bsd-determinante-schur-complemento}
 \det_F(\mathcal U_{\rm bas}^{-1}\mathscr A_E)
 =\det_F(U_{11}^{-1}A_{11})\det(\mathscr D_E).
\end{equation}

Let $\mathcal O(U_E)$ be the ring of holomorphic functions of the central disk.
The based rational spaces $V_0^i,W_0^{i+1}$ of
\eqref{eq:pdf2-bsd-s-universal} provide
\begin{equation}\label{eq:pdf2-bsd-dominios-analiticos}
 \mathscr V_E^i:=V_0^i\otimes_{\mathbb Q}\mathcal O(U_E),
 \qquad
 \mathscr W_E^{i+1}:=W_0^{i+1}\otimes_{\mathbb Q}\mathcal O(U_E).
\end{equation}
The rows of $\operatorname{Rel}_{\rm det}$ construct, over each terminal
$\eta$ and in each degree, the vectors of the common model
$v_{\eta,i}:=\upsilon_i(b_\eta)$,
$w_{\eta,i+1}:=\upsilon_{i+1}(b_\eta^{\rm term})$, and the based isomorphisms
\begin{equation}\label{eq:pdf2-bsd-jvw-generadores}
 \begin{aligned}
  J_{E,V,i}(\tau_{\eta,i}^+)&:=v_{\eta,i},\\
  J_{E,W,i+1}(\tau_{\eta,i+1}^-)&:=w_{\eta,i+1}.
 \end{aligned}
\end{equation}
The differential relation of that same row, applied to each generator, is
\begin{equation}\label{eq:pdf2-bsd-entrelazamiento-analitico-smith}
 \boxed{
 J_{E,W,i+1}(s)\,\mathscr D_E^i(s)
 =S_E^i(T_E(s))\,J_{E,V,i}(s).}
\end{equation}
This is the effective intertwiner between the corestricted Fredholm--Schur operator and the
Smith matrix; it is not a postulated equality of determinants.

To type the continuation, let $\mathfrak U_E$ be a connected domain containing
an open subset of the half-plane $\Re s>3/2$ and the disk $U_E$, and choose a finite
chain of charts $\{U_\alpha\}_{\alpha=0}^A$ of
$\operatorname{Ext}_{\Gamma,{\rm det}}$, with $U_A=U_E$. On $U_\alpha$ set
$\mathscr P_{E,\alpha}^i:=
P_{E,\alpha}^{+,i}\mathscr H_{E,\alpha}^{+,i}$ and
$\mathscr D_{E,\alpha}^i$ equal to the Schur complement of
\eqref{eq:pdf2-bsd-corestriccion-limite}. The maps to the common model are
defined, without any additional choice, by
\begin{equation}\label{eq:pdf2-bsd-j-cartas-modelo-comun}
 J_{E,\alpha,i}(\tau_{\eta,\alpha,i}):=v_{\eta,i}
 \quad\text{in the domain},
 \qquad
 J_{E,\alpha,i+1}(\tau_{\eta,\alpha,i+1}):=w_{\eta,i+1}
 \quad\text{in the codomain}.
\end{equation}
They are holomorphic because the bases of
\eqref{eq:pdf2-bsd-cambios-fredholm-cartas} are. On an overlap, define,
always in the same degree,
\begin{equation}\label{eq:pdf2-bsd-transiciones-atlas}
 G_{\alpha\beta,i}:=J_{E,\beta,i}^{-1}J_{E,\alpha,i}:
 \mathscr P_{E,\alpha}^i\big|_{U_\alpha\cap U_\beta}
 \longrightarrow
 \mathscr P_{E,\beta}^i\big|_{U_\alpha\cap U_\beta}.
\end{equation}
The literal composition of these matrices proves
\begin{equation}\label{eq:pdf2-bsd-cociclo-atlas}
 G_{\beta\gamma,i}G_{\alpha\beta,i}=G_{\alpha\gamma,i},
 \qquad G_{\alpha\alpha,i}=I,
 \qquad G_{\beta\alpha,i}=G_{\alpha\beta,i}^{-1},
\end{equation}
and the entanglement in the two letters gives the chain compatibility
\begin{equation}\label{eq:pdf2-bsd-cociclo-cadena}
 G_{\alpha\beta,i+1}\mathscr D_{E,\alpha}^i
 =\mathscr D_{E,\beta}^iG_{\alpha\beta,i}.
\end{equation}
Therefore
$g_{\alpha\beta}:=\prod_i\det(G_{\alpha\beta,i})^{(-1)^i}$ satisfies the
triple cocycle. If
\begin{equation}\label{eq:pdf2-bsd-tau-atlas}
 \tau_E^{\rm atlas}(s):=
 \prod_{\alpha=0}^{A-1}g_{\alpha,\alpha+1}(s),
\end{equation}
then $\tau_E^{\rm atlas}$ is exactly the determinantal transport
from the Euler chart to the central chart.

The even--odd totalizations $J_{E,V}$ and $J_{E,W}$ of the maps
\eqref{eq:pdf2-bsd-entrelazamiento-analitico-smith} are denoted, with
Koszul order; their determinants are
the alternating products of the determinants by degree. In the central chart
the joint unit is defined, and not before
\begin{equation}\label{eq:pdf2-bsd-u-e-construida}
 u_E(s):=\tau_E^{\rm atlas}(s)
 \det_F(U_{11}^{-1}A_{11}(s))
 \frac{\det J_{E,V}(s)}
      {\det J_{E,W}(s)}.
\end{equation}
All its factors are holomorphic and invertible upon reducing $U_E$. The
normalization is not imposed factor by factor. Let $\mathbf e_{\rm cent}$ be the
generator of the determinant line of the common model. On performing the two
triangular operations of \eqref{eq:pdf2-bsd-eliminacion-schur}, transporting
by \eqref{eq:pdf2-bsd-cambios-fredholm-cartas}, and applying
\eqref{eq:pdf2-bsd-j-cartas-modelo-comun}, the joint action on $s=1$ is,
generator by generator,
\begin{equation}\label{eq:pdf2-bsd-u-e-normalizacion-probada}
 \begin{aligned}
 &\tau_E^{\rm atlas}(1)
 \det_F(U_{11}^{-1}A_{11}(1))
 \frac{\det J_{E,V}(1)}
      {\det J_{E,W}(1)}\,\mathbf e_{\rm cent}\\
 &\qquad=
 \bigotimes_{i,\eta}
 \bigl(\epsilon_{\eta,i}^-(\tau_{\eta,i}^-)
       \epsilon_{\eta,i}^+(\tau_{\eta,i}^+)^{-1}\bigr)^{(-1)^i}
 \mathbf e_{\rm cent}
 =\mathbf e_{\rm cent},
 \qquad \boxed{u_E(1)=1}.
 \end{aligned}
\end{equation}
The first equality is the conservation identity for the common basis of the
$\mathfrak G_{\rm det}(E)$ comparator, now expanded by its rows; the
second uses
$\epsilon_{\eta,i}^\pm(\tau_{\eta,i}^\pm)=1$. None of the four factors has been
normalized separately, nor has the sought value been used.

\begin{proposition}[FERS Origin of the Analytic Factorization]
\label{prop:pdf2-bsd-fers-factorizacion-owner}
On the disk $U_E$, one has the identity of determinant sections
\begin{equation}\label{eq:pdf2-bsd-factorizacion-analitica}
 \boxed{L(E,s)=u_E(s)\det S_E(T_E(s)).}
\end{equation}
The equality is the publication of the reader FERS of
$\mathfrak G_{\rm det}(E)$ and not a equality postulada afterward of
calcular the order or the term principal.
\end{proposition}

\begin{proof}
In the absolute half-plane,
\eqref{eq:pdf2-bsd-det-a-x-l-x} identifies the Fredholm limit with $L(E,s)$.
Literal elimination \eqref{eq:pdf2-bsd-eliminacion-schur} and continuity
in trace norm give
$L(E,s)=\det_F(U_{11}^{-1}A_{11}(s))\det\mathscr D_E(s)$.
Taking ordered wedges in
\eqref{eq:pdf2-bsd-entrelazamiento-analitico-smith} gives
\[
 \det\mathscr D_E(s)
 =\frac{\det J_{E,V}(s)}
        {\det J_{E,W}(s)}
   \det S_E(T_E(s)).
\]
Equations \eqref{eq:pdf2-bsd-cociclo-atlas} and
\eqref{eq:pdf2-bsd-cociclo-cadena} transport this section through the charts;
the exact product of their transitions is $\tau_E^{\rm atlas}$. Therefore, in
$U_A=U_E$, the remaining coefficient is precisely
\eqref{eq:pdf2-bsd-u-e-construida}, proving
\eqref{eq:pdf2-bsd-factorizacion-analitica}. Uniqueness of the Mellin continuation
of \eqref{eq:pdf2-bsd-mellin-continuada} excludes another section. Finally,
\eqref{eq:pdf2-bsd-u-e-normalizacion-probada} jointly proves
$u_E(1)=1$. No step uses $d$, the rank, $\operatorname{Sha}(E)$, or the
leading term.
\end{proof}

Define, from the Smith form and not from the analytic order,
\begin{equation}\label{eq:pdf2-bsd-a-e-smith}
 a_E(T):=T^{-d}\det S_E(T).
\end{equation}
By \eqref{eq:pdf2-bsd-rboc-lt}, $a_E$ is holomorphic on $T=0$ and
\begin{equation}\label{eq:pdf2-bsd-a-e-cero}
 a_E(0)=R_{\rm Boc}^{\rm der}(S_E)\ne0.
\end{equation}
Substitution of this identity and
\eqref{eq:pdf2-bsd-t-e-construido} into the analytic factorization yields explicitly
\begin{equation}\label{eq:pdf2-bsd-expansion-local-deducida}
 \begin{aligned}
 L(E,s)
  &=(s-1)^d\,U_E^{\rm lead}(s),\\
 U_E^{\rm lead}(s)
  &:=u_E(s)
     \left(\frac{T_E(s)}{s-1}\right)^d
     a_E(T_E(s)),\\
 U_E^{\rm lead}(1)
  &=1\cdot1^d\cdot R_{\rm Boc}^{\rm der}(S_E)
   =R_{\rm Boc}^{\rm der}(S_E)\ne0.
 \end{aligned}
\end{equation}
Thus the exponent and the leading coefficient are deduced from the nonzero
holomorphic germ; they are not inputs to $T_E$, $u_E$, or $S_E$.
Now, after this deduction, define
\begin{equation}\label{eq:pdf2-bsd-dan-zan}
 d_{\rm an}:=\operatorname{ord}_{s=1}L(E,s),
 \qquad
 \mathfrak z_{\rm an}:=
 \left[(s-1)^{-d_{\rm an}}L(E,s)\right]_{s=1}.
\end{equation}
The equation \eqref{eq:pdf2-bsd-expansion-local-deducida} proves
\begin{equation}\label{eq:pdf2-bsd-dan-d}
 d_{\rm an}=d,
 \qquad
 \mathfrak z_{\rm an}
 =R_{\rm Boc}^{\rm der}(S_E)
 \quad\text{on the based line}.
\end{equation}

On the arithmetic side, the theorem~\ref{thm:pdf2-bsd-psi-pt} and the sequence
\eqref{eq:pdf2-bsd-sha-tamagawa} construct
\begin{equation}\label{eq:pdf2-bsd-dar-zar}
 \begin{aligned}
 d_{\rm ar}&:=\dim_K M_E^{\rm MW}=d,\\
 \mathfrak z_{\rm ar}&:=
 \frac{|\operatorname{Sha}(E)|\operatorname{Reg}(E)\prod_vc_v}
 {|E(\mathbb Q)_{\rm tors}|^2}
 \quad\text{on the oriented arithmetic line}.
 \end{aligned}
\end{equation}

To type the four arrows, if $C^\bullet$ is a perfect complex with
ordered basis $\mathcal B_i=(c_{i,1},\ldots,c_{i,r_i})$, one uses
\begin{equation}\label{eq:pdf2-bsd-linea-determinante-base}
 \det C^\bullet:=
 \bigotimes_i\left(\bigwedge^{r_i}C^i\right)^{(-1)^i},
 \qquad
 \mathbf e_C:=\bigotimes_i
 (c_{i,1}\wedge\cdots\wedge c_{i,r_i})^{(-1)^i}.
\end{equation}
For a finite group $A$ with based presentation
$\mathbb Z^t\xrightarrow{D_A}\mathbb Z^t\to A$, its line
$\det_{\mathbb Z}A$ carries the generator $\mathbf e_A$ and the evaluation
\begin{equation}\label{eq:pdf2-bsd-evaluacion-grupo-finito}
 \operatorname{ev}_A(\mathbf e_A):=|\det D_A|=|A|.
\end{equation}
When that line is combined with a line over $K=\mathbb Q$, the scalar
extension is always used
\begin{equation}\label{eq:pdf2-bsd-linea-finita-k}
 \det_KA:=\det_{\mathbb Z}A\otimes_{\mathbb Z}K;
\end{equation}
modules over distinct rings are not tensorized directly.
In particular, choose the Smith bases already constructed in
\eqref{eq:pdf2-bsd-smith-finito}, $m_1,\ldots,m_t$ of $M_E$, and
$n_1,\ldots,n_t$ of $N_E$, with $D_Em_i=s_in_i$, and fix
\begin{equation}\label{eq:pdf2-bsd-generador-smith-fe}
 \mathbf e_{F_E}:=
 (m_1\wedge\cdots\wedge m_t)\otimes
 (n_1\wedge\cdots\wedge n_t)^{-1},
 \qquad
 \operatorname{ev}_{F_E}(\mathbf e_{F_E})=\prod_i|s_i|.
\end{equation}
The successive lines are
\begin{equation}\label{eq:pdf2-bsd-cinco-lineas-tipadas}
\begin{aligned}
 \Delta_K&:=\det\mathcal C_E^K,\\
 \Delta_{\rm FERS}&:=\det\mathcal C_E^{\rm Iw},\\
 \Delta_{\rm PT}&:=
   \det_K M_E^{\rm MW}\otimes
   \det_K(M_E^{\rm MW})^\vee\otimes\det_KF_E,\\
 \Delta_{\rm STS}&:=
   \det_K M_E^{\rm MW}\otimes
   \det_K(M_E^{\rm MW})^\vee\otimes
   \det_K\operatorname{Sha}(E)
   \otimes\bigotimes_{v<\infty}
      \det_K\Phi_v(\mathbb F_v),\\
 \Delta_{\rm BSD}&:=\mathbb R\,\mathbf e_{\rm BSD}.
\end{aligned}
\end{equation}
The principal specialization is not a tacit identification. Denote
$I:=(T)\subset K[[T]]$ and
$\mathfrak c_{K/{\rm FERS}}^{\rm reg}:=
\det_{\rm alt}\Phi_{K/{\rm FERS}}:\Delta_K[[T]]\to
\Delta_{\rm FERS}[[T]]$. Let
\begin{equation}\label{eq:pdf2-bsd-lineas-leading}
 \begin{aligned}
  \Delta_{\rm FERS}^{(d)}&:=I^d\Delta_{\rm FERS}[[T]],&
  \Delta_{\rm FERS}^{\rm lead}
    &:=I^{-d}\Delta_{\rm FERS}^{(d)}
       \otimes_{K[[T]],\,T\mapsto0}K,\\
  \Delta_K^{(d)}&:=(\mathfrak c_{K/{\rm FERS}}^{\rm reg})^{-1}
       (\Delta_{\rm FERS}^{(d)}),&
  \Delta_K^{\rm lead}
    &:=I^{-d}\Delta_K^{(d)}
       \otimes_{K[[T]],\,T\mapsto0}K.
 \end{aligned}
\end{equation}
Here $I^{-d}\Delta^{(d)}$ denotes the module formed by the symbols
$T^{-d}z$ with $z\in\Delta^{(d)}$. The specialization is defined only
on the part of order at least $d$:
\begin{equation}\label{eq:pdf2-bsd-sp-leading-tipado}
 \operatorname{sp}^{(d)}_{\rm FERS}:
 \Delta_{\rm FERS}^{(d)}\longrightarrow
 \Delta_{\rm FERS}^{\rm lead},
 \qquad
 z(T)\longmapsto
 \overline{T^{-d}z(T)}.
\end{equation}
The Kato--FERS equivalence then induces the typed map
\begin{equation}\label{eq:pdf2-bsd-c-k-fers-leading}
 \begin{aligned}
 \mathfrak c_{K/{\rm FERS}}^{\rm lead}:{}
 &\Delta_K^{\rm lead}\xrightarrow{\ \sim\ }
   \Delta_{\rm FERS}^{\rm lead},\\
 &\overline{T^{-d}z_K(T)}\longmapsto
   \overline{T^{-d}
   \mathfrak c_{K/{\rm FERS}}^{\rm reg}(z_K(T))}.
 \end{aligned}
\end{equation}
The formula no applies $\operatorname{sp}^{(d)}$ to a linea desnuda. In
particular, if
\begin{equation}\label{eq:pdf2-bsd-seccion-zeta-leading}
 \begin{aligned}
  \boldsymbol\zeta_{\rm FERS}(T)
    &:=\det S_E(T)\,\mathbf e_{\rm FERS}
       \in\Delta_{\rm FERS}^{(d)},\\
  \boldsymbol\zeta_K(T)
    &:=(\mathfrak c_{K/{\rm FERS}}^{\rm reg})^{-1}
       \boldsymbol\zeta_{\rm FERS}(T)
       \in\Delta_K^{(d)},
 \end{aligned}
\end{equation}
then
\begin{equation}\label{eq:pdf2-bsd-especializacion-zeta}
 \operatorname{sp}^{(d)}_{\rm FERS}
       (\boldsymbol\zeta_{\rm FERS})
 =R_{\rm Boc}^{\rm der}(S_E)\,
       \mathbf e_{\rm FERS}^{\rm lead},
 \quad
 \mathfrak c_{K/{\rm FERS}}^{\rm lead}
       (\overline{T^{-d}\boldsymbol\zeta_K})
 =R_{\rm Boc}^{\rm der}(S_E)\,
       \mathbf e_{\rm FERS}^{\rm lead}.
\end{equation}
The definition uses only $d=\operatorname{ord}_T\det S_E(T)$, already computed by
Smith, and not the analytic order.

The four owner arrows, now with explicit domains and codomains,
are compared in the real fiber from the first point where the
Néron height enters. To avoid multiplying a line over $K$ by a
real scalar that does not generally belong to $K$, we fix
\begin{equation}\label{eq:pdf2-bsd-lineas-reales-comparador}
 \begin{aligned}
  \Delta_{K,\mathbb R}^{\rm lead}
    &:=\Delta_K^{\rm lead}\otimes_{K,\iota_\infty}\mathbb R,&
  \Delta_{{\rm FERS},\mathbb R}^{\rm lead}
    &:=\Delta_{\rm FERS}^{\rm lead}\otimes_{K,\iota_\infty}\mathbb R,\\
  \Delta_{{\rm PT},\mathbb R}
    &:=\Delta_{\rm PT}\otimes_{K,\iota_\infty}\mathbb R,&
  \Delta_{{\rm STS},\mathbb R}
    &:=\Delta_{\rm STS}\otimes_{K,\iota_\infty}\mathbb R.
 \end{aligned}
\end{equation}
The four arrows are then
\begin{equation}\label{eq:pdf2-bsd-cuatro-comparadores}
\begin{aligned}
 \mathfrak c_{K/{\rm FERS},\mathbb R}^{\rm lead}:
   &\Delta_{K,\mathbb R}^{\rm lead}
      \longrightarrow\Delta_{{\rm FERS},\mathbb R}^{\rm lead},
 &\mathfrak c_{K/{\rm FERS},\mathbb R}^{\rm lead}
   &:=\mathfrak c_{K/{\rm FERS}}^{\rm lead}\otimes_K1_{\mathbb R},\\
 \mathfrak c_{{\rm PT},\mathbb R}:
   &\Delta_{{\rm FERS},\mathbb R}^{\rm lead}
      \longrightarrow\Delta_{{\rm PT},\mathbb R},
 &\mathfrak c_{{\rm PT},\mathbb R}(z)&:=
   \mathfrak n_{E,\mathbb R}(z)\otimes\mathbf e_{F_E},\\
 \mathfrak c_{{\rm STS},\mathbb R}:
   &\Delta_{{\rm PT},\mathbb R}
      \longrightarrow\Delta_{{\rm STS},\mathbb R},
 &\mathfrak c_{{\rm STS},\mathbb R}&:=\det\bigl(
   0\to\operatorname{Sha}(E)\to F_E
     \to\textstyle\bigoplus_{v<\infty}\Phi_v(\mathbb F_v)
     \to0\bigr)\otimes_K1_{\mathbb R},\\
 \mathfrak c_{\rm tors,\infty}:
   &\Delta_{{\rm STS},\mathbb R}
      \longrightarrow\Delta_{\rm BSD}.&&
\end{aligned}
\end{equation}

\begin{proposition}[Generator-by-generator calculation of the four arrows]
\label{prop:pdf2-bsd-cuatro-flechas-generadores}
The actions on the bases of
\eqref{eq:pdf2-bsd-cinco-lineas-tipadas} are as follows.

\emph{First arrow.} If
$\mathbf e_K^i=\bigwedge_{b\in\mathcal B_i^K}b$ and
$\mathbf e_{\rm FERS}^i=\bigwedge_{b\in\mathcal B_i^K}\upsilon_i(b)$,
the table \eqref{eq:pdf2-bsd-tabla-kato-fers} gives
\begin{equation}\label{eq:pdf2-bsd-c1-generadores}
 \mathfrak c_{K/{\rm FERS}}^{\rm reg}(\mathbf e_K)
 =\prod_i\det(I+TB_i(T))^{(-1)^i}\mathbf e_{\rm FERS}
 =\rho_E(T)\mathbf e_{\rm FERS},
 \qquad \rho_E(0)=1.
\end{equation}
If $\mathbf e_K^{\rm lead}$ and $\mathbf e_{\rm FERS}^{\rm lead}$ are the
bases induced by \eqref{eq:pdf2-bsd-lineas-leading}, then
\begin{equation}\label{eq:pdf2-bsd-c1-leading-generador}
 \mathfrak c_{K/{\rm FERS},\mathbb R}^{\rm lead}
     (\mathbf e_K^{\rm lead}\otimes1)
 =\rho_E(0)(\mathbf e_{\rm FERS}^{\rm lead}\otimes1)
 =\mathbf e_{\rm FERS}^{\rm lead}\otimes1.
\end{equation}

\emph{Second arrow.} Choose integral bases
$a_{q,1},\ldots,a_{q,r_q}$ of $A_{q,\mathbb Z}$, viewed in $A_q$, and
$b_{q,j}=\bar\beta_q(a_{q,j})$ of $B_q$. The height matrix of the page
is defined by
\begin{equation}\label{eq:pdf2-bsd-hq-generadores}
 j_q(b_{q,j})=
 \sum_{k=1}^{r_q}(H_q)_{kj}\,a_{q,k}^\vee\otimes T^q.
\end{equation}
Therefore, on the entire wedge,
\begin{equation}\label{eq:pdf2-bsd-c2-cuna-paginas}
 \bigwedge_{j=1}^{r_q}j_q\bar\beta_q(a_{q,j})
 =T^{qr_q}\det(H_q)
  \bigwedge_{k=1}^{r_q}a_{q,k}^\vee.
\end{equation}
The product over $q$, together with the regular rows, has exponent
$\sum_q qr_q=\sum_q\dim_K V_q=d$. For
$1\leq j\leq r_q$ and $1\leq k\leq q$, set
$e_{q,j,k}:=a_{q,j}e_{q,k}$ in
$\operatorname{Jet}_{\rm Boc,\mathbb Z}(E)$ and define
\begin{equation}\label{eq:pdf2-bsd-base-mw-desde-jets}
 P_{q,j,k}:=\Psi_{{\rm PT},\mathbb Z}^{-1}
       \mathcal U_{\rm jet}(e_{q,j,k})
 \in M_{E,\mathbb Z}^{\rm MW}.
\end{equation}
By \eqref{eq:pdf2-bsd-desdoblamiento-jets},
\eqref{eq:pdf2-bsd-dimension-jet}, and
\eqref{eq:pdf2-bsd-psi-pt-integral}, these $d$ elements are an integral
basis, not a sublattice of uncontrolled index. They are enumerated in
lexicographic order as $P_1,\ldots,P_d$ and, over $K$, form the basis of
$M_E^{\rm MW}$ used below.

The height is constructed independently of FERS.\@ For
$P\in E(\overline{\mathbb Q})$, fix
\begin{equation}\label{eq:pdf2-bsd-altura-canonica-construida}
 \widehat h_E(P):=
 \frac12\lim_{n\to\infty}4^{-n}h_x([2^n]P),
 \qquad
 \langle P,Q\rangle_{\rm NT}:=
 \frac12\bigl(\widehat h_E(P+Q)-\widehat h_E(P)-\widehat h_E(Q)\bigr).
\end{equation}
Equivalently, the Poincaré biextensor rigidified over the Néron
model produces, for every place $v$, including $v=\infty$, the real local
symbol $\lambda_v(P,Q)$ normalized by the zero section, and
\begin{equation}\label{eq:pdf2-bsd-altura-neron-local-global}
 \langle P,Q\rangle_{\rm NT}
 =\sum_v\lambda_v(P,Q)\in\mathbb R.
\end{equation}
This real, symmetric, and positive height is not replaced by a sum of
cup-product invariants with values in $\mathbb Q/\mathbb Z$.

Set
$H_{\rm NT}^{\flat}(P)(Q):=\langle P,Q\rangle_{\rm NT}$ and let
$\mathcal B_E:=\Psi_{{\rm PT},\mathbb Z}^{-1}\mathcal U_{\rm jet}$,
extended to $\mathbb R$. The $q-1$ additional directions are not
replaced by identities: the full Gram block is computed, including all
cross terms among depths, rows, and jets,
\begin{equation}\label{eq:pdf2-bsd-gram-jet-completo}
 \begin{aligned}
 G_E^{\rm jet}\bigl((q,j,k),(q',j',k')\bigr)
 &:={}
 \left\langle P_{q,j,k},P_{q',j',k'}\right\rangle_{\rm NT}\\
 &=\sum_v\lambda_v
   \left(P_{q,j,k},P_{q',j',k'}\right),
 \end{aligned}
\end{equation}
for every $1\le k\le q$, $1\le k'\le q'$. This matrix is material: the
$P_{q,j,k}$ were constructed in
\eqref{eq:pdf2-bsd-base-mw-desde-jets}, and each entry is computed through
\eqref{eq:pdf2-bsd-altura-neron-local-global}. By definition of the pullback
of a bilinear form,
\begin{equation}\label{eq:pdf2-bsd-altura-nt-generadores}
 \boxed{
 \mathcal B_E^\vee H_{\rm NT}^{\flat}\mathcal B_E=G_E^{\rm jet}.}
\end{equation}
It is not asserted that $G_E^{\rm jet}$ is block diagonal. In the
integral basis $P_1,\ldots,P_d$, the regulator and its action on the determinant line
are exactly
\begin{equation}\label{eq:pdf2-bsd-regulador-cuna-altura}
 \begin{aligned}
 \operatorname{Reg}(E)&:=\det G_E^{\rm jet}>0,\\
 \bigwedge_{j=1}^dH_{\rm NT}^{\flat}(P_j)
 &=\operatorname{Reg}(E)
   (P_1^\vee\wedge\cdots\wedge P_d^\vee).
 \end{aligned}
\end{equation}
The Bockstein page and the height are two constructed pairings, with
bases already fixed. The Neron comparator between their lines is defined before
the leading term is used, by
\begin{equation}\label{eq:pdf2-bsd-comparador-neron-jet}
 \begin{aligned}
 \mathfrak n_{E,\mathbb R}:
 \Delta_{{\rm FERS},\mathbb R}^{\rm lead}&\longrightarrow
 \bigl(\det_K(M_E^{\rm MW})\otimes_K
 \det_K(M_E^{\rm MW})^\vee\bigr)\otimes_{K,\iota_\infty}\mathbb R,\\
 \mathfrak n_{E,\mathbb R}
 (\mathbf e_{\rm FERS}^{\rm lead}\otimes1)
 &:=\frac{\operatorname{Reg}(E)}{R_{\rm Boc}^{\rm der}(S_E)}
 (P_1\wedge\cdots\wedge P_d)
 \otimes(P_1^\vee\wedge\cdots\wedge P_d^\vee)\otimes1.
 \end{aligned}
\end{equation}
The denominator is nonzero by \eqref{eq:pdf2-bsd-a-e-cero}; numerator and
denominator are obtained, respectively, from the full Gram matrix
\eqref{eq:pdf2-bsd-gram-jet-completo} and from the pages
\eqref{eq:pdf2-bsd-rboc-lt}. Thus the normalization is a real quotient of two
previously computed determinants, not an invented orthogonality or an
appeal to the BSD formula. Scale extension is performed before
applying the regulator and is preserved in the PT and STS arrows.

The upper pages and the regular row preserve the integral generator of
$F_E$ by means of a typed map. The rational extension of the section
$s_q$ already constructed in \eqref{eq:pdf2-bsd-desdoblamiento-jets} is used. The dual section
\begin{equation}\label{eq:pdf2-bsd-glue-dual-tipado}
 \operatorname{gl}_q^\vee:A_q^\vee\longrightarrow
   (\operatorname{Surv}_{\rm det}^{\rm MW})^\vee
\end{equation}
is defined on the complete basis by
$(\operatorname{gl}_q^\vee\varphi)
(\operatorname{gl}_r s_r(a))=\delta_{qr}\varphi(a)$ and as zero on the
FS rows.\@ Let $\mathcal I_{\rm fin}$ be the ordered list of regular rows and pairs
$(q,j)$ with $q\geq2$. For $\eta$ with $q(\eta)\geq2$, set
\begin{equation}\label{eq:pdf2-bsd-lifts-filas-smith}
 \begin{aligned}
 \widetilde a_\eta&:=\operatorname{pr}_{M_E}\operatorname{pr}_{\mathcal L_E}
   \Theta_E\operatorname{gl}_{q(\eta)}s_{q(\eta)}(a_\eta),\\
 \widetilde b_\eta^\vee&:=
   \operatorname{pr}_{N_E^\vee}\operatorname{pr}_{\mathcal L_E^\vee}
   (\Theta_E^\vee)^{-1}\operatorname{gl}_{q(\eta)}^\vee
   \bigl((1\otimes\partial_T^{(q(\eta))})
      j_{q(\eta)}\bar\beta_{q(\eta)}(a_\eta)\bigr).
 \end{aligned}
\end{equation}
In a regular row they are defined separately
\begin{equation}\label{eq:pdf2-bsd-lifts-fila-regular}
 \widetilde a_\eta:=
 \operatorname{pr}_{M_E}\operatorname{pr}_{\mathcal L_E}
 \Theta_Es_{\rm reg}(a_\eta),
 \qquad
 \widetilde b_\eta^\vee:=
 \operatorname{pr}_{N_E^\vee}\operatorname{pr}_{\mathcal L_E^\vee}
 (\Theta_E^\vee)^{-1}s_{\rm reg}^\vee
 (\bar S_E(0)a_\eta).
\end{equation}
All domains are now $A_q$ or $A_q^\vee$ as appropriate;
$\operatorname{gl}_q$ is not applied to a covector.

The required bijection is not inferred from unimodular matrices. The
unit--counit equations show on each generator that the lists
$(\widetilde a_\eta)_\eta$ and their dual terminals are bases of $M_E$ and
$N_E^\vee$. Write the inherited list as
$\mathcal I_{\rm fin}=\{\eta_1<\cdots<\eta_t\}$ and define directly
\begin{equation}\label{eq:pdf2-bsd-nu-construida}
 \boxed{\nu(\eta_j):=j},
 \qquad
 \nu:\mathcal I_{\rm fin}\xrightarrow{\sim}\{1,\ldots,t\}.
\end{equation}
Set $\mu_j:=\widetilde a_{\eta_j}$ and let
$\xi_1,\ldots,\xi_t$ be the terminal basis of $N_E$. Before Smith, the
material matrix is
\begin{equation}\label{eq:pdf2-bsd-c2-filas-finitas}
 D^{\rm term}_{kj}:=
 \langle\xi_k^\vee,D_E\mu_j\rangle,
 \qquad
 \widetilde b_{\eta_j}^\vee
 =\sum_{k=1}^tD^{\rm term}_{kj}\xi_k^\vee.
\end{equation}
This equality is the terminal relation of
\eqref{eq:pdf2-bsd-descomposicion-global}, evaluated at the counit
$\xi_k^\vee$; it does not divide by an elementary divisor. If
$UD^{\rm term}V=\operatorname{diag}(s_1,\ldots,s_t)$, then
\begin{equation}\label{eq:pdf2-bsd-cuna-smith-con-uv}
 \bigwedge_{j=1}^t\widetilde b_{\eta_j}^\vee
 =\det(D^{\rm term})
   (\xi_1^\vee\wedge\cdots\wedge\xi_t^\vee),
 \qquad
 |\det(D^{\rm term})|=\prod_{j=1}^t|s_j|.
\end{equation}
The factors $\det U,\det V\in\{\pm1\}$ are recorded by
$\operatorname{or}_{\rm det}$ upon passing from $(\mu_j),(\xi_j)$ to the Smith
bases $(m_j),(n_j)$. Hence the wedge quotient is exactly
$\mathbf e_{F_E}$ of \eqref{eq:pdf2-bsd-generador-smith-fe}; it has not been
assumed that $U$ or $V$ are permutation matrices, and no page remains
unpublished. Consequently,
\begin{equation}\label{eq:pdf2-bsd-c2-generadores}
 \mathfrak c_{{\rm PT},\mathbb R}
 \bigl(R_{\rm Boc}^{\rm der}(S_E)
       (\mathbf e_{\rm FERS}^{\rm lead}\otimes1)\bigr)
 =\operatorname{Reg}(E)\,
  (P_1\wedge\cdots\wedge P_d)\otimes
  (P_1^\vee\wedge\cdots\wedge P_d^\vee)
  \otimes\mathbf e_{F_E}\otimes1.
\end{equation}

\emph{Third arrow.} For each $\ell$, choose ordered generators $\sigma_{\ell,a}$ of $\operatorname{Sha}(E)[\ell^\infty]$ and $\phi_{\ell,v,b}$ of $\Phi_v(\mathbb F_v)[\ell^\infty]$. By exactness, the explicit lifts $\operatorname{Lift}_{\Phi,\ell}(\phi_{\ell,v,b})$ and the images $\iota_{{\rm Sha},\ell}(\sigma_{\ell,a})$ form a presentation basis of $F_{E,\ell}$. In those bases, the relation matrix is computed generator by generator as
\begin{equation}\label{eq:pdf2-bsd-c3-matriz-presentacion}
 D_{F,\ell}=
 \begin{pmatrix}
  D_{{\rm Sha},\ell}&X_\ell\\
  0&\displaystyle\bigoplus_{v<\infty}D_{\Phi_v,\ell}
 \end{pmatrix}.
\end{equation}
The column $X_\ell$ records the change of each lift when its
representative is modified; the lower-left block is zero because
$\partial_{{\rm loc},\ell}\iota_{{\rm Sha},\ell}=0$. Therefore the presentation wedge
does not depend on $X_\ell$ and
$\det D_{F,\ell}=\det D_{{\rm Sha},\ell}
\prod_{v<\infty}\det D_{\Phi_v,\ell}$. On it,
\begin{equation}\label{eq:pdf2-bsd-c3-generadores}
 \mathbf e_{F_{E,\ell}}
 \longmapsto
 \mathbf e_{\operatorname{Sha}(E)[\ell^\infty]}
 \otimes\bigotimes_{v<\infty}
    \mathbf e_{\Phi_v(\mathbb F_v)[\ell^\infty]},
\end{equation}
and the evaluations of \eqref{eq:pdf2-bsd-evaluacion-grupo-finito} give
$|F_{E,\ell}|=|\operatorname{Sha}(E)[\ell^\infty]|
\prod_{v<\infty}|\Phi_v(\mathbb F_v)[\ell^\infty]|$.
Multiplying over $\ell$ produces
\begin{equation}\label{eq:pdf2-bsd-c3-integral-generadores}
 \mathfrak c_{{\rm STS},\mathbb R}(\mathbf e_{F_E}\otimes1)
 =(\mathbf e_{\operatorname{Sha}(E)}
  \otimes\bigotimes_{v<\infty}\mathbf e_{\Phi_v(\mathbb F_v)})\otimes1,
 \qquad
 |F_E|=|\operatorname{Sha}(E)|\prod_vc_v.
\end{equation}

\emph{Fourth arrow.} Given the preceding Mordell--Weil basis and a
Néron differential $\omega_E$ with positive period
$\Omega_E$, the oriented basis of $\Delta_{\rm STS}$ is sent by
\begin{equation}\label{eq:pdf2-bsd-c4-generadores}
 \begin{aligned}
 &x\,(P_1\wedge\cdots\wedge P_d)\otimes
 (P_1^\vee\wedge\cdots\wedge P_d^\vee)\otimes
 \mathbf e_{\operatorname{Sha}(E)}
 \otimes\bigotimes_{v<\infty}\mathbf e_{\Phi_v(\mathbb F_v)}\\
 &\hspace{28mm}\longmapsto
 \frac{x\,\Omega_E\,
   \operatorname{ev}_{\operatorname{Sha}(E)}
      (\mathbf e_{\operatorname{Sha}(E)})
   \prod_{v<\infty}\operatorname{ev}_{\Phi_v(\mathbb F_v)}
      (\mathbf e_{\Phi_v(\mathbb F_v)})}
 {|E(\mathbb Q)_{\rm tors}|^2}\,
 \mathbf e_{\rm BSD}.
\end{aligned}
\end{equation}
By \eqref{eq:pdf2-bsd-evaluacion-grupo-finito}, the coefficient of the last
line is
$x\Omega_E|\operatorname{Sha}(E)|\prod_vc_v/
|E(\mathbb Q)_{\rm tors}|^2$.
The two copies of the torsion group follow from the dual degrees $0$ and $2$;
the factor $\Omega_E$ is the evaluation of $\omega_E$ on the oriented real
component. This orientation places $\Omega_E\mathfrak z_{\rm ar}$, and not its
inverse, as the arithmetic coefficient of the composition; the comparison with
$\mathfrak z_{\rm an}$ will be deduced only subsequently through the square of
trivializations.
\end{proposition}

\begin{proof}
The first identity is the alternating determinant of the triangular matrices $I+TB_i(T)$. The second applies the specialization of the pages and then the Neron comparator \eqref{eq:pdf2-bsd-comparador-neron-jet}. Its numerator is the determinant of the complete jet Gram matrix \eqref{eq:pdf2-bsd-gram-jet-completo}, with all cross-terms, and its denominator is the Bockstein determinant already calculated. By \eqref{eq:pdf2-bsd-regulador-cuna-altura}, its action on wedges is \eqref{eq:pdf2-bsd-c2-generadores}; the exponent $d$ is obtained before taking Mordell--Weil dimensions. The third is the canonical isomorphism of lines associated with the exact sequence, written in the already constructed sections $\iota_{{\rm Sha},\ell}$ and $\operatorname{Lift}_{\Phi,\ell}$. The fourth is the evaluation of the two finite torsion complexes and of the Neron line. This calculates the four arrows without using the BSD equality.
\end{proof}

Its composition
\begin{equation}\label{eq:pdf2-bsd-comparador-total}
 \mathfrak C_E:=
 \mathfrak c_{\rm tors,\infty}\circ
 \mathfrak c_{{\rm STS},\mathbb R}\circ
 \mathfrak c_{{\rm PT},\mathbb R}\circ
 \mathfrak c_{K/{\rm FERS},\mathbb R}^{\rm lead}:
 \Delta_{K,\mathbb R}^{\rm lead}
 \longrightarrow\Delta_{\rm BSD}
\end{equation}
is well typed. Before inserting any analytic or cardinal element,
the independent trivializations are fixed
\begin{equation}\label{eq:pdf2-bsd-trivializaciones-independientes}
 \operatorname{tr}_{\rm an}(x\mathbf e_K^{\rm lead}):=x,
 \qquad
 \operatorname{tr}_{\rm Ner}(y\mathbf e_{\rm BSD}):=y.
\end{equation}
The first uses only the Kato basis specialized by
\eqref{eq:pdf2-bsd-lineas-leading} and the Mellin--Fredholm normalization of
\eqref{eq:pdf2-bsd-fredholm-euler}; in particular,
$\operatorname{tr}_{\rm an}
(\overline{T^{-d}\boldsymbol\zeta_K})
=R_{\rm Boc}^{\rm der}(S_E)=\mathfrak z_{\rm an}$ by
\eqref{eq:pdf2-bsd-especializacion-zeta}. The second uses only the oriented Neron
differential and satisfies
$\operatorname{tr}_{\rm Ner}(\mathbf e_{\rm BSD})=1$. The identity of the based
comparator of $\mathfrak G_{\rm det}(E)$, unfolded by the four
arrows, is the
square
\begin{equation}\label{eq:pdf2-bsd-cuadrado-trivializaciones}
\begin{array}{ccc}
 \Delta_{K,\mathbb R}^{\rm lead}
   &\xrightarrow{\ \mathfrak C_E\ }&\Delta_{\rm BSD}\\[2mm]
 \big\downarrow\scriptstyle\operatorname{tr}_{\rm an}
   &&\big\downarrow\scriptstyle\operatorname{tr}_{\rm Ner}\\[2mm]
 \mathbb R&=&\mathbb R.
\end{array}
\end{equation}
This square is fixed with the structural bases before evaluating
$L(E,s)$, the regulator, or the finite groups.

The analytic element is placed, after constructing the arrows, in its
correct domain:
\begin{equation}\label{eq:pdf2-bsd-elemento-analitico-tipado}
 \mathbf z_E^{\rm an}:=\mathfrak z_{\rm an}\mathbf e_K^{\rm lead}
 =\overline{T^{-d}\boldsymbol\zeta_K}
 \in\Delta_{K,\mathbb R}^{\rm lead}.
\end{equation}
By \eqref{eq:pdf2-bsd-dan-d},
$\mathfrak z_{\rm an}=R_{\rm Boc}^{\rm der}(S_E)$; the second equality in
\eqref{eq:pdf2-bsd-especializacion-zeta} proves that the first arrow carries
$\mathbf z_E^{\rm an}$ exactly to the left-hand side of
\eqref{eq:pdf2-bsd-c2-generadores}. The other three arrows, computed on
generators, give
\begin{equation}\label{eq:pdf2-bsd-elementos-linea}
 d_{\rm an}=d=d_{\rm ar},
 \qquad
 \boxed{
 \mathfrak C_E(\mathbf z_E^{\rm an})
 =\Omega_E\mathfrak z_{\rm ar}\,\mathbf e_{\rm BSD}.}
\end{equation}
Applying the vertical arrows of
\eqref{eq:pdf2-bsd-cuadrado-trivializaciones} produces the numerical equality
\begin{equation}\label{eq:pdf2-bsd-igualdad-numerica-deducida}
 \boxed{\mathfrak z_{\rm an}=\Omega_E\mathfrak z_{\rm ar}.}
\end{equation}

\begin{theorem}[Birch--Swinnerton--Dyer Publication]
\label{thm:pdf2-bsd-publicacion-final}
For the complete autonomous determinant image
$\mathfrak G_{\rm det}(E)$, one has
\begin{equation}\label{eq:pdf2-bsd-rango-final}
 \operatorname{ord}_{s=1}L(E,s)=\operatorname{rank}E(\mathbb Q)=:r
\end{equation}
and
\begin{equation}\label{eq:pdf2-bsd-formula-final}
 \frac{L^{(r)}(E,1)}{r!\,\Omega_E}
 =\frac{|\operatorname{Sha}(E)|\operatorname{Reg}(E)\prod_vc_v}
 {|E(\mathbb Q)_{\rm tors}|^2}.
\end{equation}
\end{theorem}

\begin{proof}
The equality of degrees is
\eqref{eq:pdf2-bsd-dan-d} followed by
\eqref{eq:pdf2-bsd-rango-orden}. The Kato--FERS table transports the analytic element
to every Bockstein page; $\Psi_{\rm PT}$ glues them to the
Mordell--Weil lattice and publishes its height determinant; the exact
Sha--Tamagawa sequence decomposes the Smith factor; and the final comparator inserts
the torsional lattice and the already oriented period. This is exactly
\eqref{eq:pdf2-bsd-elementos-linea}. The independent square
\eqref{eq:pdf2-bsd-cuadrado-trivializaciones} gives
\eqref{eq:pdf2-bsd-igualdad-numerica-deducida}; substituting the definitions of
$\mathfrak z_{\rm an}$ and $\mathfrak z_{\rm ar}$ yields
\eqref{eq:pdf2-bsd-formula-final}. None of the four arrows is
normalized with $r$, $\operatorname{Reg}(E)$,
$|\operatorname{Sha}(E)|$, or the leading term: those values appear
only at the end of their respective determinants.
\end{proof}

\section{Non-governing controls, falsifiers and status}

\begin{falsifier}[Falsifiers separated by component]
The construction fails if a history of
$\Sigma_{\rm det}^{\rm mult}$ does not belong to the partition
\eqref{eq:pdf2-bsd-tres-proyectores}; if one writes
$\partial H_{\rm FS}+H_{\rm FS}\partial=q_{\rm root}$ and thereby erases
$\mathcal L_E$; if a nonroot row of
\eqref{eq:pdf2-bsd-tabla-kato-fers} has zero weight; if the unit--counit equations
\eqref{eq:pdf2-bsd-unidad-counidad-paginas} fail; if
$D_E\otimes\mathbb Q$ is not invertible; if
$\pi_\Phi\operatorname{Lift}_\Phi\ne I$; or if any arrow of
\eqref{eq:pdf2-bsd-cuatro-comparadores} uses the final value to fix its
base. Each falsifier breaks a specific edge and does not reopen the preceding ones.
\end{falsifier}

The homotopy $H_{\rm FS}$ is a subsequent control of the explicit filler
$h_*=-vf$ and not a governing premise. Smith forms,
determinants and cardinalities are computed after constructing the maps;
they do not replace the tables, the gluing of pages or the exactness of cocycles.

\paragraph{Status and provenance.}
Alexander--Whitney, the fiber product, the common unit, reduced
duality, root rigidity, the filler, all Bockstein pages and
the four readers belong to
\texttt{PREEXISTING\_AUTHORIAL\_ARCHITECTURE}/\texttt{INTERNALLY\_EXACT} of
$\mathfrak G_{\rm det}(E)$. The exhaustive partition
\eqref{eq:pdf2-bsd-descomposicion-global}, the table
\eqref{eq:pdf2-bsd-tabla-kato-fers}, the formulas
\eqref{eq:pdf2-bsd-psi-y-lambda-q} and the cocycle maps
\eqref{eq:pdf2-bsd-iota-sha-cociclos}--
\eqref{eq:pdf2-bsd-lift-componentes} are a
\texttt{NEW\_FORMALIZATION}: they materially unfold the intertwiners
already contained in the thirteen structural coordinates. The final theorem has the
force
\texttt{PROVED\_WITH\_AN\_EXPLICIT\_STARTING\_STRUCTURE}; no
governing residual obligation remains, nor is an abstract homotopy used to
manufacture the result.
